\documentclass[letterpaper]{article}
\usepackage{url}
\usepackage{natbib}
\usepackage{multicol}
\usepackage{fullpage}
\usepackage{mdwlist}
\usepackage{graphicx}
\usepackage{framed}
\usepackage{listings}

\newcommand{\serc}{S$^2$ERC}

\title{What a Cybersecurity Education Game Reveals About How Children Conceive of Security and Software Engineering\\
\medskip
\serc{} Technical Report \textbf{NNN}\\
}

\author{Paul Gestwicki and Kaleb Stumbaugh\\
Computer Science Department\\
Ball State University\\
Muncie, IN 47306\\
Email: pvgestwicki@bsu.edu}

\begin{document}

\maketitle


\begin{abstract}
  Abstract goes here.
\end{abstract}

\section{Introduction}

The Cybersecurity Education Workshop Final
Report~\cite{Cybersecurity2014} identified six major research themes
for immediate action, in order to address deficiencies in contemporary
computer security education practice. The first of these---\textit{Concepts and Conceptual Understanding}---was the topic of our earlier work
on cybersecurity epistemology, particularly as applied to
educational games~\citep{Gestwicki2015,Gestwicki2015-tr}.
Our analysis of the literature and a survey of the state of the art
contributed to a three-tier taxonomy for cybersecurity education games:

\begin{description}
\item[Type 1] Games that convey cybersecurity concepts through
 narrative and/or theme only. There is no representation of
 the concepts within actual gameplay. That is, the act of playing
 the game does not require any decision-making that would reflect
 an understanding of cybersecurity concepts.
\item[Type 2] Games that integrate multiple-choice questions
  (including yes/no options and branching narratives)
  that correspond to cybersecurity concepts.
  Answering these prompts correctly requires an understanding
  of the concepts.
\item[Type 3] Games that require ambiguous decision-making
 such that making good decisions implies an understanding of
 cybersecurity concepts.
 \end{description}

In this report, we turn to two additional themes from the
Cybersecurity Education Workshop Final Report:
\textit{Assessment} and \textit{Recruitment and Retention}.
The former deals with the development of reliable and valid
techniques for measuring subjects' understanding of cybersecurity
as well as methods for evaluating educational interventions.
The latter deals with the ``pipeline'' problem, that there are not enough
people preparing for and staying in cybersecurity careers to meet the demand.

This work focuses on youth in middle school and early high school. Children at
this age make critical decisions about their futures, particularly with regard
to whether or not they can succeed as scientists 
and engineers. \citet{Margolis2003} and \citet{Margolis2010} describe
how cultural factors disproportionately affect underrepresented groups,
contributing to their absence in IT careers generally.

We embarked on an iterative and incremental
approach to game design and development, following best practices
of agile game software development.
The result of this process is \textit{The Social Startup Game}, 
a Type~3 cybersecurity education game that explores themes of
technology careers, educational paths, and security and
software engineering fundamentals.
In particular, we designed and evaluated a game designed around
two learning outcomes: that the player would learn fundamentals of
cybersecurity, and that the player would learn about careers and
career paths in cybersecurity.
Section~\ref{sec:design} of this report documents our design process
and results. Section~\ref{sec:evaluation} describes our
qualitative evaluation protocol and results, and Section~\ref{sec:discussion}
contextualizes these results within the overarching research framework.


\section{Game Design}
\label{sec:design}

\subsection{Early Multiplayer Prototypes}

Following the analysis presented in \citet{Gestwicki2015}, 
we proceeded with a rapid prototyping process in order to identify
an optimal design for our constraints.
We focused on physical prototypes
as we worked through the various cybersecurity epistemologies,
matching concepts to appropriate game mechanisms.

The MDA framework~\citep{Hunicke2004} was useful here, 
as it separates the \textit{mechanics} (the elements that a designer
controls) from the \textit{dynamics} (emergent gameplay properties of the
mechanics). This is analogous to many security vulnerabilities, in which
the system designer specifies the rules of a system (mechanics) and the
attacker attempts to exploit unexpected combinations and applications
of that system (dynamics).

Also useful for cybersecurity game development is the observation
that real-life security is a competitive, asynchronous, distributed,
asymmetric game. That is, security is:
\begin{description}
\item[competitive] attackers versus defenders
\item[asynchronous]  the two sides may take their actions at any time,
 regardless of what the other is doing
\item[distributed] the two teams are not collocated, although
 allies may be
\item[asymmetric] the two sides take different actions
\end{description}
These properties create an interesting and significant design challenge.
Like all creative and scientific endeavors,  game design generally
derives new ideas from previous ones, punctuated by occasional
paradigmatic shifts and inventions. In our analysis, we found
no games with these properties. The closest we encountered were
digital capture-the-flag games, which are timeboxed and not strictly
asynchronous.

The \textit{White Hat / Black Hat} (WHBH) family of designs
represent attempts to convey the competitive and asymmetric properties
of cybersecurity to the player.
The first of these designs featured an asymmetric design in which
one player was the Black Hat and the other 2--4 players were 
White Hats. White Hats needed to secure company assets, while the
Black Hat simultaneously tried to capitalize on others' successes.
The main tension 
tension in this design came from the White Hat players desire
to beat each other and outsmart the Black Hat, while the Black Hat
player was able to hide exploits and sow mistrust.

A new WHBH design emerged from this process, one that cast each player
as both a White Hat and a Black Hat: each round of gameplay consisted
of a White Hat phase in which players managed company assets, followed
by a Black Hat phase in which players secretly sabotaged each other.
More specifically, a White Hat hacker has to balance the business
demands of their company against the need for confidentiality,
integrity, and availability of their services, while a Black Hat
hacker seeks to find exploits and has to balance the risk of discovery
against the increased rewards that come with patience.  The main
conflict in this game, then, came from managing the unknowns: is your
system compromised? Will you be attacked? Will your opponent discover
your virus before you can gain valuable data?

Asymmetry was kept as a focus by having the White Hat and Black Hat
phases play differently. The White Hat phase included fairly conventional
resource management gameplay.
The Black Hat phase involved mini-games---puzzles that represented,
at an abstract level, authentic hacking activities.
For example, a tile-rotation puzzle was designed to replicate
finding a path through a system and the inherent delay between
writing a program and verifying its execution.

In parallel with the WHBH family of designs, we developed and evaluated 
a playable prototype of an untitled multiplayer resource management game.
This game also had players acting as both white hat and black hat hackers,
but this game featured more explicit representation of four
important values: confidentiality, integrity, availability, and profit.
White hat players could either reinvest their profits directly in
profit-making activity, or they could strengthen their defenses.
During the black hat turn, players played secret attacks against 
each other, requiring a third party (not one of the players) to maintain
some of the game's secret state.
Playtesting this prototype, we learned that 
while the concepts of confidentiality, integrity, and availability
resonate with security professionals who already understand what
they represent, these terms---as abstractions---do not carry much meaning
to those not in IT. 
The game design treated these terms separately and authentically,
each being a defense against a different kind of attack,
but this was arbitrary to the players.

\subsection{Genesis of a single-player simulation game}

After approximately two months of rapid prototyping, we decided
to abandon multiplayer designs and move to a simpler single-player
design. The two authors had previously worked on a realtime
multiplayer video 
game\footnote{\url{http://collaborationstationgame.info}} and
were wary about the complex production needs for such designs, especially
given the amount of other design risk being encumbered.

We were quickly able to take what we had learned from the rapid
prototyping experience and transform this into the seed of a single-player
simulation game.
Although a single-player design lost some authenticity in representing
the adversarial nature of cybersecurity, it also eliminated an underlying
conundrum we had faced during prototyping: would potential audiences
such as teachers, librarians, or parents 
reject a game that forced a young
player to simulate illegal (black hat) activities?

Considering the importance of choosing metaphors familiar to the player---as
opposed to terms like ``integrity'' that are abstruse to our target
audience---we adopted a theme of social media.
Informal polling showed that youth in our target demographic 
were familiar with social media, particularly through popular brands
such as Facebook and Instagram. 
We also decided to keep the balance represented in our previous
multiplayer designs, where security concerns are ``defensive'' moves
that are necessary but not sufficient for victory.

One of the drawbacks of our earlier, multiplayer designs was that they
cast each player abstractly as both a person and a company, and this
elided the fact that software development is cooperative activity.
We were inspired by the cooperative game theory, which observes
that software development is a cooperative game of invention
and communication~\citep{Cockburn2006}.
This led to our decision to frame the player as a security expert
with a team of developers working under them.
This theme allowed us to abstract away money as a directly manipulable
resource and instead to focus on social media's fungible currency: users.

\subsection{\textit{Social Startup Game}}

\textit{Social Startup Game} is a single-player simulation game
in which the player is cast as the chief security advisor for
an up-and-coming social media startup, Social 
Jam;
a list of the game's key concepts is provided in
Appendix~\ref{appendix:key-concepts}.
The player manages three employees over ten days and is given two 
competing goals: first, to increase the number of Social Jam users
to 20,000, and second, to keep the estimated exposure (security risk)
of the company at less than 15\%.
The player does this by assigning employees to either the Development
task or Maintenance task. Development contributes to the completion of
features, which earn new users but also increase exposure.
Social Jam is constantly under threat by hackers, and the higher the
exposure, the higher the chance of an exploit. The Maintenance task
is necessary for both fixing exploits and reducing exposure.

There are nine
main characters (listed in Appendix~\ref{appendix:characters}), and at
the start of the game, one is randomly chosen to be CEO, and three
more for your three employees.
Embedded into the simulation game are the eight narrative events
listed in Figure~\ref{fig:event-list}; these events serve as
important decision points for contextualizing educational content.
%
\begin{figure}
\begin{framed}
\begin{description}
\item[Introduction] Explains the game interface.
\item[Take-your-child-to-work Day] A coworker asks the player to 
 recommend a major for his daughter to pursue in college.
\item[Data stolen] Choose responsible disclosure or secrecy when 
 data is stolen by hackers.
\item[DDOS] Respond to an angry ex-user's initiation of a distributed
 denial-of-service attack.
\item[Insecure password] When a weak password causes a problem,
 choose whether to ignore it, fix it, or train the staff to prevent
 future problems.
\item[Script kiddie] Faced with amateurish attacks, choose to 
 ignore it, report it, or fight back.
\item[Security conference] One employee may be sent to a two-day
 conference to improve their Maintenance Skill
\item[Input sanitization] When an employee expresses concern over
 input sanitization, you can either take action or tell them to ignore it.
\end{description}
\caption{Narrative events of \textit{Social Startup Game}}
\label{fig:event-list}
\end{framed}
\end{figure}
%
One event is the introduction that explains the game's interface,
and the sequencing and timing of the other eight
are randomized at the start of the game. 

At the end of ten in-game days, the player is given feedback from
the Social Jam CEO about their performance. This feedback varies
depending on how well the player met the two different goals, the best
feedback of course being reserved for those who met both.

\subsection{Technical Details}

\textit{Social Startup Game} is a Web-based game based on HTML5 and
runs without plugins in all modern evergreen browsers. The aspect
ratio of the game makes it friendly for mobile devices such as smart
phones and tablets, and the game scales to any screen size.  The game
is hosted at \url{socialstartupgame.info}; it is open source, and the
central repository is at
\url{github.com/doctor-g/social-startup-game}.

The game was written in Java using PlayN~(\url{playn.io}),
which integrates with GWT~(\url{www.gwtproject.org})
to compile the Java source code into
HTML5 and JavaScript. 
\textit{Social Startup Game} is a Maven project, and other
libraries directly 
used by \textit{Social Startup Game} include: 
\begin{itemize}
\item TriplePlay (\url{github.com/threerings/tripleplay}), 
 a library of utilities for PlayN-based games, including 
 screen and UI management
\item Mockito (\url{mockito.org}), a mocking framework for unit tests
\item Google Guava (\url{github.com/google/guava}), a collection
 of useful libraries for Java
\item React (\url{github.com/threerings/react}), a library for
 functional-reactive programming in Java
\end{itemize}

The game has a hybrid architecture, combining elements of traditional
object-oriented design, an entity system, and functional-reactive
programming.
Most of the simulation is implemented within the entity system,
which is a style of game engine architecture that offers many benefits over
standard OO conventions~\citep{Gestwicki2012}. The entity system
approach divides the software into \textit{entities} that
aggregate \textit{components}, and \textit{systems} that process 
specific kinds of components.
For example, there is an entity that represents a feature, and that
entity has a user-generating component. The user generation system
looks at all active entities that have such components and updates
the Social Jam user count based on this.
\textit{Social Startup Game}'s entity system comprises 15~systems
and 25~components, the explanations for which are provided
in Appendix~\ref{appendix:entity-system}.

All of the programming and design were done by Paul Gestwicki and
Kaleb Stumbaugh, and all of the artwork was created by Coy Yuan.  The
music is by Kevin MacLeod (\url{incompetech.com}) and licensed under a
Creative Commons license (CC-BY-3.0).


\section{Evaluation}
\label{sec:evaluation}

\subsection{Methods}

To answer our research questions requires building an understanding of
the player from descriptions and observations, and so we follow a
qualitative methodology~\citep[see][for example]{Stake2010}.  The
paucity of cybersecurity educational game assessments leaves almost no
grounds for forming the reliable and valid instruments necessary for
quantitative methods. Put another way, we could not know whether what
we measured would actually contribute to answering our research
questions.  \citeauthor{Cohen1994}'s classic essay highlights several
methodological problems with applying statistical methods to human
behavior, not the least of which is the common confusion of
\textit{probability of data} with \textit{probability of hypothesis
  correctness}~\citep{Cohen1994}.
Qualitative research methods, by contrast, deal with 
observations and descriptions, meeting the subjects within their
complex cultural contexts, and seeking understanding of phenomena
that are not directly measurable. 

Our research protocol was approved by Ball State University
Institutional Review Board~(835177-1). 
Our subject pool consisted of thirteen youth, ages ten to fifteen,
with eight females and five males. All of the subjects were
white American students living in Indiana. 
They came from families who responded to a call for participation that
was sent to Ball State University mailing lists and shared via social
media. This group is a convenience sample of those who responded to the
call, chosen for their ability to schedule sessions and with an emphasis
on female participants. The subjects include a mixture of homeschool
and public school attendees.

\begin{table}
\begin{framed}
\centering
\begin{tabular}{llr}
\textbf{Subject ID} & \textbf{Gender} & \textbf{Age}\\
1 & Female & 14\\
2 & Male & 10\\
3 & Female & 12\\
4 & Female & 11\\
5 & Male & 12\\
6 & Male & 15\\
7 & Female & 12\\
8 & Male & \textbf{??}\\
9 & Female & \textbf{??}\\
10 & Female & \textbf{??}\\
11 & Male & 15\\
12 & Female & 10\\
13 & Female & 10\\
\end{tabular}
\caption{Research Subjects}
\end{framed}
\end{table}

Each session consisted of a semi-structured interview combined with
playtesting~(Appendix~\ref{appendix:interview-protocol}.
Interview audio was digitally recorded, and these were transcribed
into 918 paragraph-separated units. 
The investigators also 
produced~27 paragraphs of analytic memos,
% PVG 1140, KCS 229
consisting of~1378 words.
Gameplay behavior was automatically tracked using integrated logging, 
producing~16 logs consisting of~326 discrete game events.

Interview data were coded using techniques inspired by
\citet{Saldana2009}. The two researchers conducted three phases 
of coding before comparing results in order to ensure inter-rater
reliability. The first round consisted of open coding, and these
codes were refined in the second round. 
Appendix~\ref{appendix:phase-2} lists our second phase codes,
consolidated and standardized between researchers.
The third round consisted of
identifying categories, which represented the concepts that emerged
out of multiple second-round codes; the categories are listed in
Appendix~\ref{appendix:categories}.

\subsection{Findings}

Our analysis resulted in four major themes, which are listed
in Table~\ref{tab:themes}.
%
\begin{table}
\begin{framed}
\begin{itemize}
\item Mixed views on education and career goals
\item Diverse opinions about developers appearance and interests
\item Impact of background knowledge
\item Two modes of character-driven decisions
\end{itemize}
\caption{Four Major Themes}
\label{tab:themes}
\end{framed}
\end{table}
%
Each of the themes identified by our analysis are described in turn below.


\newcommand{\theme}[1]{\subsubsection*{#1}}

\theme{Mixed views on the role of education and degree toward career goals}

The subjects articulated a predictable mix of opinions about their
school experience. Although the interview question about favorite
school subjects was intended as more of an ice-breaker than anything,
many students eagerly described what subjects they enjoyed and why.
For example, subject~7 enjoyed English and science because they allow her
to be ``more creative than in other classes,'' and subject~5 likes
mathematics because identifies as ``a problem-solver.''
Subject~6 gives his favorite subject as lunch, since he is ``not a big
fan of school\ldots{} I know I have to do it, but it's not very fun.''
We observe some subjects making distinctions between ``academic'' and 
``non-academic'' courses: subject~9 prefers art and band, but she identifies
English has her favorite academic course. We lack data to indicate whether
this distinction is from schooling structures, environmental impacts (her
mother is a teacher), or deeper knowledge of epistemology.

Our subjects showed a general understanding that mathematics is important
to being a developer. According to subject~11, 
``There's probably a lot of mathematics and, you know, how you have to
know how games work or what other people did to make the games.''
% 
They also perceived Computer Science as being a default path into
application development.  Subject~6 says, ``It seems like making an
app would be more computers and that stuff than anything else.''
During the in-game event in which players have to choose what major to recommend
to a coworker's daughter, ten chose Computer Science, two chose 
Graphic Design, and one chose Engineering---a surprising consistency
of choices given that there were nine options (Table~\ref{tab:child-advice}).
%
\begin{table}
\begin{framed}
\centering
\begin{tabular}{rlrl}
\textbf{Option} & \textbf{Times Chosen}\\
Computer Science & 10\\
Engineering & 1\\
English & 0\\
Graphic Design & 2\\
Mathematics & 0\\
Marketing & 0\\
Performing Arts & 0\\
Physics & 0\\
Psychology & 0
\end{tabular}
\caption{Possible college major choices to suggest to a coworker's daughter 
 in \textit{Social Startup Game} and the number of times they were
 selected by unique subjects during interviews}
\label{tab:child-advice}
\end{framed}
\end{table}
%
Although almost all the students refer to Computer Science as the
subject to study to become a developer, there is no other data
to suggest that they understand what Computer Science is beyond a gateway
to application development.

Some of the subjects recognized that a college education is not the only
path to becoming a successful application developer.
Subject~5 identifies reputation as an asset and entrepreneurship as a 
possibility, saying that ``\ldots{} somebody would hear about you, and you probably need to get a degree, or be an entrepreneur and start your own app-making business.''
Subject~6 recognizes that coursework can help someone learn to make apps
but that there is a wealth of ``random sources'' from which one can learn
as well:
\begin{quote}
If somebody decides they want to make 
an app, and they know how, they sit down and do it. If they don't know how,
and they still want to do it, then depending on how complex the app is they 
either go to try to figure it out at some random source, or if it's a bigger,
more complex project, they go take classes for it, and then they make the app.
\end{quote}

Subject~4 recognized that teachers hold significant sway in students'
opportunities. When asked how developers got t heir jobs, she suggested
that ``their teachers saw their ability to do something
like that in an earlier stage, so when they were able to get
the classes in high school and through college to be able to get better
at that.''
This empowerment of authority is an interesting contrast
against the entrepreneurial view described above,
yet it reflects a real and common understanding of the trajectory
from elementary school through to a career.



\theme{Diverse opinions about developers' appearances and interests}

Most of the players claimed to have never thought about the physical
characteristics of software developers, but among those that did,
there was great variance.  Subjects 6~and 13 reported developers to be
``geeky,'' which matched our expectations based on earlier work on
perceptions of engineers (\citet{Margolis2003}, for example). Neither
subject elaborated on what characterized geekiness, and we recognize
that there are potential misunderstandings in leaving the terms used
by teenagers up to the interpretation of adults.  
By contrast,
subject~10 stated that app developers look ``like normal
people\ldots{} like anybody.''
Subject~10 went on to say that the characters in the game
``didn't look like smart people---they looked like normal people.''

We did not expect the responses that painted app developers as something
attractive and aspirational. Subject~7 describes them as ``17 to sort of
early 20s in age,'' only about five to ten years older than herself.
Subject~9 was amused by Ivar's eccentric appearance, saying, ``He looks
kinda creepy, and I love it. I like his creepy long hair with a bald head.''
This was particularly amusing to us as Computer Scientists, given that
Ivar was modeled after Bjarne Stoustrup, but we did not bring this up
in the interview.

Subject~4 brings up a distinction between what she calls ``designers''
and ``producers,'' which we believe to correspond to the business functions
of application or interface design and programming. 
When asked what app developers look like, she puts it like this:
\begin{quote}
Well, I'm pretty sure [they look] like people, they're like artists
and stuff, \ldots{} I picture the people who \textit{design} it would
be kind of like artists, kind of, messy kind of regular people, and the
people who publish it would be kind of like, \textit{official}.
\end{quote}
It is worth noting that subject~4 was one of several students who had
completed a business information technology course in the local public
school, which we know to have included an introduction to programming
and which we believe to have informed students about some
specializations in the field.

Subject~6 does not like school, but he does
like \textit{Minecraft}. When asked about the appearance of \textit{Minecraft}'s
developers, he responds, ``Probably most of them are not American,
probably a lot of them are from Japan. It's a Japanese company, isn't it?''
While other subjects who brought up \textit{Minecraft} seemed to know something
about its Swedish creator, Markus Persson (``notch''), subject~6 held
an incorrect belief that it was created by Japanese developers. 
We suspect that this belief may be due to the relatively high number
of Japanese game development companies, this does not come up in the interview.

There was only one instance of overt sexism in the data:
when asked why he sent Esteban to the security conference over
Janine or Melissa, subject~2 responded cheekily, ``He's the guy.''
No other subject referenced developers' gender in their discussion
of their perceptions, backgrounds, interests, or educational preparation.
However, we recognize that gender and racial prejudices may be
subconscious or intentionally hidden from subjects' responses.



\theme{Impact of background knowledge}

We saw players' background knowledge impact their in-game behavior
and gameplay outcomes in a variety of ways.
In a theoretical sense, this is unsurprising, given the wealth of
evidence to support constructivism---the educational theory that
individuals build their own mental models based on their experiences
and background knowledge.
However, the broad and multifaceted impact that this can have on
educational game design and evaluation makes this worth analyzing.

Most players expressed surprise at the scale and persistence
of cyberattacks. They were less clear on what the hackers
were after. Subject~2 expressed a concern that revealed a few
layers of confusion about the game, saying that perhaps the
hackers ``wanted to call the game their own.''
This shows an understanding of software copyright infringement (``piracy''),
which then shaped his rationalization for their behavior.
Although some subjects showed an understanding of the business and
economic forces around software development, none of our subjects
reflected aloud on the value of private information.

In the hacker event, most of our subjects chose to report the
crime to the police. Subject~5, however, had a different
perspective, saying this while playing the game: 
\begin{quote}
I could report it to the police, or I could task someone to try to
figure it out. I'll probably just task somebody to figure it out for now
because, I don't really feel like getting mixed up with the police,
so I'll task [someone] to do it.
\end{quote}
We did not follow up to inquire why he did not want to get ``mixed up
with the police,'' but it stands in sharp contrast to those subjects
who, seeing this option, quickly decided it was the best course of
action.  Subject~10 also chose to seek vengeance over
notifying the police---an action that results in a tremendous loss of
users---and in her post-game reflection, she regretted it:
``No, I probably shouldn't have. I probably should have let
the police handle it instead of taking it into my own hands.''

Several subjects thought that decision points had right and wrong
answers, and that their task was to choose the right one.
This matches their experience with education and conventional
educational games. Given that their interaction with us was framed in terms
of an educational game, it is not surprising that this mental model
would be at the fore.
Subject~11's explanation of how he approached decision-making in the
game could just as well be about how he took a quiz in school:
``I knew some of them would hurt me no matter what I did so [I chose] the
best possible answers. And some of them, I really didn't know what to do,
so I just guessed.''
However, we also observed some players adopting a more nuanced view,
particularly those who played more than once. Subject~9 puts it this way:
\begin{quote}
See, last time I was like, ``Oh, I'll let them know,'' and everyone
got mad. So I thought I'd just keep it a secret.
But then when I was exposed, everyone was like,
 ``Why didn't you tell us that?!''
They all got really mad and left me, which I think is what made my exposure
go up. Last time they didn't like that; this time, they didn't like it more.
\end{quote}

The players did not question whether the game was an authentic
portrayal of a social media company, although they recognized
the game as being a simulation.
Subject~5 says, ``I actually kind of liked it. I liked the way you have to try and run your company---I just kind of like that type of thing.''

Finally, we note that students' vocabulary and cultural literacy
were major factors in their ability to understand the gameplay
experience.
None of the players younger than~14 understood the term ``press release,''
although they were able to understand their options when we explained
the term.
Subject~5 demonstrated a general misunderstanding of core
gameplay terms when, a few minutes into the game, he asked,
``Are exploits and features good or bad, what are these?''
Game literacy also played an important role, as it appears many of
our subjects did not recognize the character skill levels for what they were.
By contrast, upper-division high school players had no trouble recognizing
these, and we can assume that older players would have broader game literacy.


\theme{Different modes of character empathy}

The players often referenced their workers when making decisions. Several
took advantage of the game's paused state to review the bios of their
employees. We observed a dichotomy in how this information was used.
Those who engaged in \textit{pragmatic} decision-making chose
characters who were deemed to be best at these roles. This is exemplified
in subject~4's approach to dealing with the hacker event:
``They needed somebody to fix it, so
I sent [Vani] because she was good at fixing that kind
of thing.'' Similarly, subject~7 explains why sent Jerry to investigate
the injection attack, contrasting his background against Janine's:
``He seemed like the most well-trained out of all of them,
and the girl---I forget her name---she seemed more like someone
who should be with, like, social, and not developer,
but Jerry seemed like the better person for the job of that kind of thing.''
Sometimes this pragmatic approach included assumptions about 
educational background as well, as we can see in subject~11's 
explanation of why he referenced their bios so often:
\begin{quote}
I thought there was probably a pretty balanced team, because you had
someone who majored in three of the important categories,
such as computer engineering.
I can't remember what the one in the middle was, and
then there was another one who was good at working with people and
finding out what they want. That's why I chose here to do the---oh no,
I chose the engineering one to teach people how to do the password thing because
he know how to make a safe password, and then I would have chosen her to
find out what people liked\ldots{}.
\end{quote}

The other mode of character-based decision-making was \textit{empathic}:
some players justified their decisions based less on a desire to
maximize output and, instead, based on a feeling of helping the
fictional characters feel more fulfilled.
Subject~9 is the best example of this---an enthusiastic player
who quickly learned the characters' names, referenced them frequently,
and seemed particularly interested in being a good leader to them.
When we asked why she chose Janine for several events, she told us, 
``From her description, she seemed really on it. She seemed really cool.
I liked her a lot.'' 
Subject~9 also commented on how much she liked Nancy, explaining that
she saw some of herself in Nancy's artwork:
\begin{quote}
She looked really sassy. And she looked like me before
I went insane, 'cause I used to have long blond hair and that was
normal. It was gross. But she was really sassy, so I liked her.
`Sassy,' like, not really sassy but, like, `cool.' Like, she looked hip
and with the kids.'' 
\end{quote}
It is worth adding that, when subject~9 was first reading Nancy's 
educational credentials, she had other reasons to identify with
this character: Nancy and the subject's mother shared an \textit{alma mater},
and she was also impressed with Nancy's hobbies, saying, ``Nancy has a popular podcast about being a woman in technology. Oh, that's cool!'' 
Our data are not clear on whether sharing a gender contributed to this
identification as well, but we note the possibility.


\section{Discussion}
\label{sec:discussion}

\subsection{About the game}

Our findings regarding game vocabulary affirm our design decision
to avoid the terms ``confidentiality'', ``integrity'', and ``availability''
in the game. These are undeniably security jargon, not common parlance.
The subjects in our study were not a random sample of their age group,
and in fact, we have reason to believe that these were children of
privilege; we expect that players' confusion over vocabulary
and cultural references would be exaggerated among players from
less privileged backgrounds.

\textit{The Social Startup Game} features high degrees of randomization
in the spirit of increasing replayability. For example, the sequence
of story events and the set of characters included in the game are randomized,
so they will be unpredictable for any given gameplay experience.
However, we confess that while this improves replayability---generally
considered a virtue in games of any kind---it also potentially hindered
our evaluation. Each player was faced with a random set of characters,
for example, so we cannot directly compare two players' opinions of their
workers to each other, and we also could not intentionally match
or mismatch character gender and ethnicity with the player.

A significant category in our analysis included usability problems
with the game interface.
The two that dominated the list were players not realizing that they
could change characters' tasks and that they could
read character bios. This surprised us during the first few interviews,
given that all of the players read this introduction given by Frieda,
Social Jam's administrative coordinator:
\begin{quote} 
  [Employee 1], [Employee 2], and [Employee 3] are currently maintaining our
  software. You can tap them to find out more about them. You may
  reassign any number of them to new feature development at any
  time. Go ahead and try that now, and let me know when you are ready!
\end{quote}
All of the players read this and then tapped ``Next'' without ever following her
instructions, except for those whom we interrupted to demonstrate
how to do these actions.
This usability problem could likely have been discovered with 
more playtesting prior to the formal evaluation. However, we do not
believe that this had any invalidating effect upon our findings,
since players were quick to understand the interface once we
demonstrated how to use it. We suspect part of the confusion came from
the fact that our interface was designed for touch-based mobile devices,
yet our evaluation was being done on non-touch laptops with mouse
input.\footnote{The reason for using mouse-driven input was
 an untimely defect in one of our supporting libraries. 
PlayN's 2.0-rc3 contains a previously unidentified defect 
that broke popup menu support on touch-input devices, and hence player role
selection was impossible on touch-based interfaces. Although PlayN
2.0 has yet to be released, the snapshot build has fixed this defect,
and our current version of the game is built upon that snapshot.}
We observe with mouse-driven input that players scan the screen
with the pointer, watching for responses that would indicate 
interactivity, whereas on touch-based devices, players are more
apt to try tapping things to see what happens.

We used \textit{The Social Startup Game} as part of our workshop 
at the 2016 Congressional Leadership Academy, which is a full-day event
for high-achieving high school juniors from Indiana's sixth district.
The students at this event, who were just barely outside of our
target demographic, showed none of the confusion over terms
and jargon; in fact, they seemed to enjoy the game much more than our
target demographic. 
Of course, these were high-achieving students who elected
to attend a session on cybersecurity, so in a sense we were preaching
to the choir: these students were already on their way to potential
careers in cybersecurity. 


\subsection{About security and software development}

As described in Section~\ref{sec:design}, our game veered away
from the fine-grained security concepts explored in 
games such as \textit{CyberCIEGE}~\citep{Irvine2005} toward
authentic contextualization of relevant decision-making.
We saw variation in how subjects approached their initial choices,
informed by their individual backgrounds, but we also saw a move toward
more ethical and informed decision-making after playing the game.
This was especially pronounced when considering the principle
of responsible disclosure.
In the absence of a longitudinal component to our evaluation,
we do not know whether these students will later recall specific
kinds of attacks introduced in \textit{The Social Startup Game},
such as injection and DDOS attacks.

Our players who had previous exposure to programming believed that
they could get a job in software development. This is an important
victory for advocates of early computing education intervention.
One homeschooled subject mentioned having a positive experience with
Code.org, while several public-school subjects cited positive experience
with Scratch\footnote{\url{https://scratch.mit.edu}} through a required
business information technology course. Subjects with no exposure to
programming showed more anxiety or unease at the thought of software
and IT careers. It merits repeating here that ours was
not a random sample: our subjects came from families who were willing
to bring them to an educational game evaluation, and all our subjects
identify as ``white.'' Fortunately, there are ample other programs---even
within our own community---that are building local theories around
ethnic minorities and specifically targeting families of low socio-economic 
status.
Despite this caveat around ethnicity, we observed no significant difference
in opinions or articulations between our male and female participants.
We hope that this points to a continuation of the trend toward greater
gender diversity in Computer Science programs and, by extension, into
industry.


\subsection{About education}

The number and confidence of players who chose Computer Science
as the default academic option for IT careers is problematic.
While it is true that many students graduate from Computer Science
programs do pursue such careers, we know that there are many other
ways to get there. Anecdotally, we have met recruiters who prefer
liberally-educated, critical-thinking philosophy majors over 
myopic, code-oriented Computer Science majors.
If there is any standard of what constitutes Computer Science,
it is the ACM/IEEE Curriculum Guide~\citep{CS2013}---a sprawling
guide that speaks to the truth of the insider joke, ``If you ask
five computer scientists to define `computer science' you get
seven different answers.''
We are starting to see new programs being developed, including
Software Engineering programs and transdisciplinary programs like
Georgia Tech's computational media major. However, the fact remains
that to many students, computer science has ``computer'' in the title,
and it looks to them like the one surefire path toward software
development careers.

We believe there is reason to be concerned about the perception that
subjects in school are either ``academic'' or ``not academic.''
Our study does not reveal from whence these youth developed
such perspectives, although we can assume that it relates to the
administrative structure of their school experience.
The divisions of school are artificial, created by bureaucracy for
the purpose of perceived efficiency.
Students who do not recognize this as an artificial construction 
become, ironically, victims of it rather than liberally educated.

Although this study is aimed at youth, our experiences with
undergraduate computer science students revealed to us a connection
between our subjects and our university students: 
understanding of security principles and practices is negligible, 
and both pools seem to lack the concepts---the vocabulary---about
which to converse on these. This echoes the findings of the 
Cybersecurity Education Workshop~\citep{Cybersecurity2014}.
We suspect that there is a causal relationship between the
prevailing structures of Computer Science education and the lack
of good security practices in industry. The ACM/IEEE curriculum
guidelines for Computer Science curricula treat cybersecurity differently
from all other domains, explaining that this topic must be included
throughout the curriculum and not simply in isolated courses.
However, conventional higher educational structures privilege the
separation of content into courses as well as the academic freedom
of faculty in the classroom. This leads to a situation where there is
no practical way to create interleaving topics, motivate their adoption,
or evaluate their efficacy. 


\subsection{About game design}

Those who made pragmatic character-based decisions did so assuming
that character backgrounds impacted their performance, while in fact,
all of the characters are interchangeable: the differences between
them are only skin deep. 
It is interesting that the same players who lacked game literacy to
recognize certain design tropes (such as a character statistics block)
did assume that character backgrounds provided some kind of hint to
maximizing decision outcomes.
Considering two design options---one in which characters are all fundamentally
the same regardless of their appearance, bio-sketch, and credentials, and one where their in-game behavior is different based on these factors---we lack data
to determine whether one will better meet our intended learning outcomes.

The discovery of empathic character-based decisions opens more
potential opportunities for future design research.
Our prototyping process described in Section~ref{sec:design} included
several different quantitative factors such as users, defects, and money,
but none of our approaches considered emotional health, such as happiness
or job satisfaction. Players who tend toward empathic decision-making
may feel more engaged with the game by receiving such feedback, and this,
in turn, may lead to better learning outcomes.
Inspired by Bartle's classic essay~\citep{Bartle1996} in which he
describes the kinds of players who play MUDs, perhaps a similar
kind of taxonomy could be built that describes how players approach
decision-making with respect to fictional, in-game characters.
We know of no such theory, but our findings suggest that it could be
quite useful for game design generally and educational game design
specifically.

Many players made decisions based on a search for the ``right'' answer,
not just a good answer. We suspect that there are two factors
that contribute to this mental model. First, the game was framed as
an ``educational game'', which is something of a loaded genre.
\citet{Klopfer2009} describe the history of educational games and their
many problems. One common design mode is to rely upon multiple-choice,
quiz-style assessments, where there is just one right answer.
We explored how this phenomenon affects cybersecurity games in particular
in our earlier work~\citep{Gestwicki2015}. 
Such assessments ignore constructivist
and situated theories of learning (see \citet{Hickey2003}, for example).
More to our point, they put students into a school-culture mental model
that multiple-choice scenarios imply that there is only one correct
answer.
The second factor we suspect to contribute here is the predominance
of such dichotomous decision trees in computer games. From our experience,
games aimed at youth tend not to feature nuanced decision-making, particularly
in dialog options. Hence, a player whose mindset is one of conventional
games would not expect, say, that the advice given to a coworker's
daughter did not have any ``wrong'' answers.

As mentioned earlier, we had the opportunity to introduce several
high school juniors to \textit{Social Startup Game} through the 
Congressional Leadership Academy. While this was not part of our
formal evaluation, we noticed an interesting phenomenon that is worth
noting here. This portion of the academy was held in a computer science
teaching lab that is designed for pair programming, and so based on
the furniture configuration, the participants formed pairs around
workstations. Each play through the game was therefore completed
by a pair of participants---strangers who had only met that day.
Although they were not instructed to follow a think-aloud protocol
as were our actual research subjects, they played the game much
slower and more deliberately, seeking consensus among the two players
before making any move. We believe that this points to an important
possible direction for future evaluations of such pro-social games.
This kind of situative perspective---drawing upon the tradition
of \citet{Wenger1999}---could reveal important truths
about the subjects and their play experience~\citep{Hickey2003}.


\section{Conclusions and Future Work}

The design goals of this project were to create a game that teaches
fundamentals of cybersecurity and to expose careers and career paths
to cybersecurity careers.
Regarding the former, we found that the topic of 
computer security poses significant challenges to the design of
educational interventions, particularly because security is woven
into the tapestry of professionalism in software development.
While it's true that cybersecurity can be a focused area of study---as
evidenced by courses and certificate programs---to understand security
requires an operational understanding of the rest of the system.
This property of cybersecurity is the reason why the ACM/IEEE Computer Science
Curriculum guidelines uniquely identify information assurance as needing
to be integrated through other courses; this is especially surprising given
the conventional, course-structured approach for the rest of the curriculum
guide.

Despite these challenges, our rapid prototyping process identified
several possible mappings from cybersecurity concepts to learning game
mechanisms. 
A limited production budget and schedule contributed significantly to the
selection of ideas that constitute \textit{Social Startup Game}.
Future work should more rigorously explore the asynchronous, asymmetric,
multiplayer mechanisms that were out of scope for this project,
since these properties are integral to computer security---we suspect
they represent the best effort to ``find the game in the 
content''~\citep{Klopfer2009}.
The context of deployment and evaluation is also worth investigating,
particularly extending to  more situated environments,
adopting more wholeheartedly a design-based
 research approach~\citep{Brown1992,Barab2004}.

Regarding the second design goal, our evaluation has served to 
confirm the value of early computing interventions toward positive
outlook of youth toward computing careers.
Our subjects recognized the economic viability of computing careers.
However, this was coupled with a wide variety of perspectives
of the developers themselves: some positive and aspirational, some
negative and stereotypical.
Taking \textit{Social Startup Game} as a pilot project, future work
could look more carefully at how both aesthetic and formal properties
of in-game characters can change perspectives of players toward
real-life software developers and IT careers.
Our findings present a snapshot, a local theory on how youth perceive
software development careers, and future work can compare this to 
other times, places, and contexts.

This work validates our hypothesis that a Type~3 cybersecurity
education game would be useful for addressing our design
goals~\citep{Gestwicki2015}. 
\textit{Social Startup Game} puts the player into an authentic 
although simplified simulation of a software development start-up company.
Players made decisions based on what was best for the company, 
aligning in-game victory with successful decision-making.
They identified with the in-game characters in rich ways, making
both pragmatic and empathic choices in their attempt to win.
This produced a corpus of data that led to important, interesting, and
sometimes surprising findings---findings that we believe would not have
been possible in a Type~1 or Type~2 experience.
This work therefore provides an example of and a justification for
formal evaluation of cybersecurity education games. 


\section{Acknowledgments}

This research is based upon work supported by the National Science
Foundation under Grant No.~0968959 and by the Ball State University
Honors College.


\bibliographystyle{plainnat}
\bibliography{references}

\clearpage
\appendix

\clearpage
\section{\textit{Social Startup Game}: Key Concepts}
\label{appendix:key-concepts}

\begin{itemize}
\item \textit{Player}: the person playing this single-player game

\item \textit{Company}: Social Jam, the company employing the player
 \begin{itemize}
 \item \textit{Exposure}: the probability of a successful exploit
 \item \textit{Users}: the number of Social Jam users, which functions as
                       the player's score
 \item \textit{Acquisition Rate}: The number of users generated per hour
 \end{itemize}

\item \textit{Employee}: A non-player who works for Social Jam under 
                        the player's direction
 \begin{itemize}
 \item \textit{Development Skill}: Productivity at development task
 \item \textit{Maintenance Skill}: Productivity at maintenance task
 \item \textit{Background}: Educational credentials and interests
 \end{itemize}

\item \textit{Task}: One of the jobs to which an employee can be assigned
 \begin{itemize}
 \item \textit{Development}: Advance completion of features
 \item \textit{Maintenance}: Lower company exposure and contribute to fixing
                             of exploits
 \item \textit{Special}: Special tasks are assigned through decisions in 
                         narrative events and have a variety of effects
 \end{itemize}

\item \textit{Exploit}: A named defect that continuously drains a small
                        percentage of users until it is fixed

\item \textit{Feature}: A user-facing feature offered by Social Jam 
                        that attracts new users but also increases
                        exposure

\item \textit{Narrative Event}: An interactive narrative moment
                         in which the player must make a choice

\end{itemize}

\clearpage
\section{Characters in \textit{Social Startup Game}}
\label{appendix:characters}

\newenvironment{character}[3]
{
\begin{tabular}{p{1.25in}p{4in}}
\raisebox{-0.5\height}{\includegraphics[height=2in]{images/#1.png}} &
\textbf{#1 #2}
\begin{itemize*}
#3
\end{itemize*}
}
{
\end{tabular}
\bigskip
}

\begin{character}{Esteban}{Cortez}
{\item Bachelors in Computer Science, Ball State University}
Esteban worked in a factory until he was 33, then he went to college and decided to get involved in software development.
\end{character}

\begin{character}{Nancy}{Stevens}
{\item Bachelors in English, Georgetown University
 \item Masters in Computer Security, Purdue University}
Nancy has a popular podcast about being a woman in technology.
\end{character}

\begin{character}{Jerry}{Chen}
{\item Bachelors in Computer Science, University of Hong Kong}
Jerry interned at a local company in high school and has been working as a software developer ever since.
\end{character}

\begin{character}{Vani}{Mishra}
{\item Bachelors in Computer Engineering, Indian Institute of Science
 \item Masters in Software Engineering, Ball State University}
Vani was born in India and came to the United States for graduate school. She loves music, dancing, and PHP.
\end{character}

\begin{character}{Abdullah}{Nasr}
{\item Bachelors in Electrical Engineering, Iowa State University}
Abdullah used to work for a larger social media company, but he prefers the excitement of a small startup.
\end{character}

\begin{character}{Janine}{Palmer}
{\item Bachelors in Computer Science, Virginia Tech}
Janine is especially talented at meeting with customers and understanding what they want from a product.
\end{character}

\begin{character}{Melissa}{Barnard}
{\item Bachelors in Mathematics, Stanford University}
Melissa has always loved games and puzzles, but she especially loves bicycling and walking her German Shepherd.
\end{character}

\begin{character}{Ivar}{Johansen}
{\item Bachelors in Chemistry, Stockholm University
 \item Masters in Electrical Engineering, Uppsala University}
Ivar teaches kids how to build simple robots as a volunteer in a local school.
\end{character}

\begin{character}{Bruce}{Powers}
{\item Bachelors in Computer Science, Ball State University}
Bruce was recently married and enjoys going to the gym and watching documentaries about history.
\end{character}

\begin{tabular}{p{1.25in}p{4in}}
\includegraphics[width=0.6in]{images/Frieda.png}
&
\textbf{Frieda}

Frieda is the administrative coordinator at Social Jam.
\end{tabular}

\clearpage
\section{\textit{Social Startup Game} Entity System}
\label{appendix:entity-system}

\subsection*{Systems}
\begin{description}
\item[EventTriggerSystem]
 Checks the registered events and triggers and event to occur if its registered time has passed
\item[ExploitMaintenanceSystem] Produces progress towards ``fixing'' the oldest exploit based on the skill of the employee(s) assigned to the Maintenance task
\item[ExposureReductionSystem] Reduces the company’s exposure value based on the maintenance skill of the employees assigned to the “maintenance” task
\item[FeatureDevelopmentSystem] Increases the progress towards completion on current feature based on the “development” skill of the employee(s) assigned to the “development” skill
\item[FeatureGenerationSystem] Generates a new feature with zero progress when the previous feature is completed
\item[GameTimeSystem] Moves forward in real time but can be paused and unpaused; used by any other system that relies on time.
\item[LayerPositionsSystem] Animates sprite layers
\item[LearningSystem] Increases an employee’s skill over time as they perform a related task
\item[TaskProgressSystem] Calculates the time remaining on special event tasks assigned to employees
\item[UpdatingSystem] Updates entities; used as a model level support of updating UI components
\item[UserAttritionSystem] Removes users per section per exploit equal to the product of the total company users and the exploit attrition value (0.0022)
\item[UserGenerationSystem] Generates users per second equal to the sum of the user generation values of the completed features
\item[WorkHoursSystem] Limits the work hours of employees to the common work day; accelerates time to twice as fast when not during normal work hours
\end{description}



\clearpage
\section{Semi-structured interview protocol}
\label{appendix:interview-protocol}

\subsection*{Before playing the game}

\begin{enumerate}
\item What is your favorite subject in school?
\item What is your favorite app?
\item Who are the people who make [favorite app]?
\begin{enumerate}
\item What do they look like?
\item Did they go to college? What did they study?
\item What were their favorite subjects in school?
\end{enumerate}
\end{enumerate}

\subsection*{After playing the game}

\begin{enumerate}
\item Did anything surprise you?
\item What kinds of decisions did you make in the game? How did you make those decisions?
\item Where do apps come from?
\item Do you think you could get a job making apps? Why or why not?
\item How does someone get a job making apps?
\item Do you have any questions?
\end{enumerate}

\clearpage
\section{Second phase codes}
\label{appendix:phase-2}
\begin{multicols}{2}
\begin{verbatim}
academics: computer science
academics: science-engineering
academics: dislike

apps: communication
apps: games
apps: productivity
apps: for phones
apps: creative
apps: music

characters: pragmatic
characters: empathic

culture: male domination
culture: smart people look different

developers-like: math
developers-like: english
developers-like: art + stem
developers-like: technology
developers-like: science
developers-like: art
developers-like: design
developers-like: programming
developers-like: animation
developers-like: fundamentals
developers-like: computer science
developers-like: engineering

developers: not considered
developers: smart
developers: work hard
developers: crazy

developer-appearance: artists
developer-appearance: nonconformist
developer-appearance: geeky
developer-appearance: young
developer-appearance: smart and techy
developer-appearance: normal

development: difficulty
development: scale
development: coding
development: teams
development: testing
development: time
development: design separate from programming
development: entrepreneurship
development: anybody can do it

distribution: distinction of sources
distribution: companies as gatekeepers
distribution: appstore

economy: demand for apps
economy: jobs

education: degree has value
education: degree not necessary
education: degree as career
education: development classes
education: online resources
education: multiple paths

experience: scratch
experience: school
experience: code.org
experience: gamesalad

self-efficacy: technology requisite

game design: role-playing
game design: one right answer
game design: genre preferences
game design: empathy
game design: strategy

interest: problem-solving
interest: Minecraft
interest: creativity
interest: english
interest: science
interest: math
interest: writing
interest: terraria
interest: art
interest: band
interest: sketchbook pro
interest: history
interest: facebook
interest: wattpad
interest: gym
interest: clash of clans

learning: working hard
learning: improvement
learning: feedback loops

ssg: usability
ssg: vocabulary
ssg: math

security: disclosure
security: hackers
security: piracy
security: responsibility
security: pragmatism
security: encryption
security: confidentiality
security: number of exploits
security: strong passwords
\end{verbatim}
\end{multicols}

\clearpage
\section{Categories}
\label{appendix:categories}
\begin{itemize}
\item[] Computer science as the domain for app development
\item[] Character-based decisions
\item[] Academic vs non-academic learning
\item[] Incremental vs Entity theory of intelligence
\item[] Varied perspectives on developer appearance
\item[] Development is difficult
\item[] Companies as gatekeepers
\item[] Democracy of development
\item[] Economics and development
\item[] Degree as career
\item[] Positive educational experiences with coding
\item[] There is one right answer
\item[] Learning through practice
\item[] SSG issues
\item[] Surprise at hackers
\item[] Responsible disclosure
\item[] Protecting personal information
\item[] ``App'' means ``game''
\end{itemize}

\clearpage
\section{Interview Transcripts}
The following pages contain transcripts from our semi-structured
interviews. The transcriptions include both the interviewer's
and the subject's statements, usually interleaved. Vertical space
indicates a pause, such as when playing the game.

\lstset{
basicstyle=\small\ttfamily,
columns=flexible,
breaklines=true
}

\section*{Subject 1}

\begin{lstlisting}
What is your favorite subject in school?
Um..probably language arts
What's your favorite app?
Like, an app that you get or just one that's on your phone?
Any
Um, well, messaging, and then, that's about it
Who do you usually talk to when messaging---school friends?
Yeah
Remind me how old you are
Fourteen
[offtopic]
Tell me about how you imagine the people who make the messaging app that you like. Where do you think those come from?
Um [looks perplexed]
Let me start with a different question: when you think about the people who make those apps, what do they look like?
Um, kind of, anything, I guess
Do you ever think about what they might have studied in school or in college?
Technology things
What do you think their favorite subject in school would be?
Maybe like, I don't know, some kind of ... math, or language arts too
You started by saying math, and you added language arts, and to a lot of people those things are very different.
Yeah
Can you tell me why you think both of those?
Well, with like, the way it's designed, and the way you can communicate with other people
OK.
[Introduce the game]
Do you have any questions ..?
At the beginning it's a little hard to tell what to do.

Did anythign surprise you?
Um, not really
What kinds of decisions did you have to make?
Whether to kind of notify people about what happened or keep it secret.
And how did you make those decisions?
Um, I guess I just kind of went with what I thought would be a good idea to do if i was in that position.
How do you know if it's a good idea or not?
Hm [shrug] cuz I don't have that kind of experience.
So what do you think about where apps come from?
I know that it usually takes a lot of people just to design little things, cuz in school in BIT last year we worked with a program called Scratch, and it was a lot harder than we thought to do stuff, to make like a bigger app or even like a little game
Can you tell me what BIT stands for?
Yeah I think it's Business In Technology but I'm not sure
Do you think you could get a job making apps?
[inaudible, but affirmative]
Lately that's what the world has been focusing on like social media, and apps and devices, and I guess there's just a lot of job opportunities for people to do that. 
So how do you think somebody who makes apps gets that job?
Um, not sure, maybe would have some kind of like start-up program where you can show what you want to make to a company or what you want to do.

[Show her how to change tasks to employees]

Can you tell me about how you're making those decisions, about who's doing what?
Um, trying to make them equal for both sides [development & maintenance skill]

Can you tell me about why everyone is on maintenance now?
Um, cuz there's more of the, I'm not sure what that one is (dev't or maintenance skill)

[game ends]
Was anything different that time?
It was a little different because I didn't know I could switch between the maintenance and development
Did that affect how you made decisions at all?
Not really.
\end{lstlisting}

\clearpage\section*{Subject 2}
\begin{lstlisting}
You're homeschooled, right?
Yeah.
Good, my kids are home schooled too. I don't know what kind of curriculum you follow, but can you tell me what your favorite kind of work is or what your favorite subject is...?
Science mostly
Science? What kinds of stuff do you do in science?
Experiments, learn about ... science, pretty much. Um... [asks mother] What science do we do Mom?
You answer!
I forget!
That's OK, tell me, how old are you?
I'm ten
Ten great. 
[off topic]
So do you use any apps?
I use... for school? or just...
At all. anything, and do you have any favorites?
I like to make pretend videos, like when I'm playing games on my mom's iPad, so yeah.
How does that work?
I just act like I'm talking to the people, and I just say what I'm doing and how I'm doing it and what I'm going to be doing and how I'm going to do it.
And you record it on the iPad then?
Yeah.
Neat. So do you ever think about the people who make those apps? They had to come from somewhere, right?
Sometimes
What do you think they look like?
I don't know, really
Try to imagine, someone who makes those apps, can you picture what they might look like or where they might be?
Not really.
OK, can you think maybe what their favorite subjects in school were?
Huh. Maybe like Math, because you kind of have to know math to kind of develop the games and do all the coding and stuff.
OK, anything else?
Not really.
Do you know any coding?
A little bit. 
What have you done?
Not really much. I've just like typed into my computer and like sometimes--most of the time nothing happened.
That's what happens to me too, and that's my job. It's hard!
[introduce game]
[Intro Ends]
So ...
[unintelligible] I just kind of...
[Shows how to click on the Alerts button]
[Demonstrate how to change task]

I can change all these guys too?
You bet

Can you tell me about what you're thinking while you're reading?
[reading screen] One of your engineers is 'connected' about, what does that say?
Injection
Injection attacks. If it's not han...dled handled correctly hackers may be able to run their programs on your machine. Who will you ... assign... to ... intestigate.

What is our goal? [Points to status screen] "Progress towards goal"?
Um, I don't remember what the number was that Frieda said. I think it's 20,000 users. Or, it was the CEO at the beginning. I think it was 20,000 users. Progress tells you what percent you are. Have you done percents in math?
Mmm mmm, not yet.
 Oh, so that might be a little confusing

Do you know what a press release is?
No.
That's when you make a public statement about what happened, you tell everybody, "Here's what happened."
[Selects it]

So the goal is so many users?
Yes, you can see your percent is over 100 now, so you've met your goal, and you have a couple hours left before the end of the game---in-game hours

[end game]
Did anything in the game surprsie you?
Like how there was _hackers_, how there was so many of them, and they were so big of a... so big of a problem to the game and stuff
And that was surprising to you?
A little bit
What do you think about the hackers?
...I don't really know.
Do you have any idea why they were attacking?
Maybe they wanted to call the game their own or something?
What kinds of decisions did you have to make in the game?
Like to see if we should tell the public or just keep it quiet about the hackers and... stuff.
I noticed when you had the option to send someone to the security conference, you sent Esteban.
Mmm hmm
Was there a reason you sent Esetban?
He's the guy [a bit cheekily stated]
When you had to investigate... I forget which attack it was, script kiddie attack maybe ... that time though you chose Janine to investigate the problems. Do you know why you chose Janine for that one?
No.
So, after playing the game do you have any thoughts about where apps come from?
Little bit.
Can you tell me about what you think that might be like?
Um,... it's like the people have to make them .. and then they have to... I'm confused. So, the apps, where they come from?
mmm hmm!
I don't really know, actually.
Do you think you could get a job making an app?
Probably not
Why is that?
I'm not that smart at making stuff when it comes to like games and stuff.
OK. [offtopic] Do you have any thoughts about how somebody  would get a job making apps?
If they're like really smart and they know like how computers work and how devices work so they know what's the best idea and stuff.
Do you have any questions for me, about the game or the study?
Not really.
\end{lstlisting}

\clearpage\section*{Subject 3}
\begin{lstlisting}
How old are you?
Twelve.
[offtopic]
What is your favorite subject when you're doing schoolwork?
probably reading
Why is that?
Just because I like the books that have a lot of, like, problems, but then they find out what the problem is and fix the problem.
Great. What is your favorite app?
I don't really have a favorite app but I do like Minecraft a lot.
What do you like about Minecraft?
I can build houses and like other buildings and stuff
Have you ever seen any pictures of notch, the guy who made minecraft.
If you imagine... Minecraft was created by a guy named notch, that's his nickname, what's his real name? Marcus Persson I think is his real name. but, now it's worked on by many many people. Can you imagine the people who make minecraft, what they might look like?
Mmm no not really
Can you imagine what kinds of things they would have liked in school?
Science, math [inaudible]
Why do you think science and math
Math because you need to know math to do like, stuff, for the game, and science because it will help you put stuff together [raises hands like slotting pieces together] to control the game, I guess
That's a good answer, that's fine. The people who made minecraft, do you think they went to college and can you imagine what they might have studied in college?
No, I don't think so.
It's a long ways away when you're twelve.
Yes
So, do you want to try our game?
Sure
[introduce the game]

OK, so, how do I figure out what to do?
Start by clicking on the blinking alerts button

do I press this? (employee area)
Sure.

[Explain the task assignment]
So I can like, press on there? (task assignment)
Mmm Hmm

[Brief talking over each other, confirming that she can use the task assignment buttons on other characters too]

What is 'maintenance'?
THat's a good question. Maintenance is keeping the current systems up and running, making sure everything is working OK, and then development is making new features.

OK, now, do I press this again?
MmmHmm

OK, so, ...
What were you going to say?
I was just going to ask you what I should do next, but ... should I change their jobs again?
You're welcome to do that, mm hmm

What does 'press release' mean? [Reads it emphasizing verb meaning of 'press']
A press release is when the company goes to the public and says, here's what's currently happening, they would tell everybody that there's a problem.

Ummm

Can you tell me about why you're making those changes?
Just to see how they work better on the... [inaudible, probably implying 'other tasks']

What does 'development' mean again?
Development is working on new features for the program

What does 'train staff' mean? [in the password event]
So you could teach your staff to use better passwords.
[selects it]

[ game ends]

Did anything surprise you in the game?
"Mm, probably how like the hackers and like the passwords and stuff. How there are hackers trying to get people's information, and how we lost people just because of [them?]
[Note that 'hackers' was a similar answer to 368.2, whose interview she was in the room for]"
Can you say anything more about why that surprised you?
mmmm I don't really know like...
It's OK. What kinds of decisions did you have to make in the game?
Figure out which people would be the best people to fix the problem and like doing something about it, not just ignoring it.
And how did you make those decisions? How did you decide who to send?
Just by what their jobs were, like 
What do you mean, when you opened them up and looked at their backgrounds?
Yeah
I noticed that when you had an option to send someone to the security confernece, you sent Nancy. Do you remember why you chose Nancy?
No, not really
And when you were looking at the injection attacks, you looked over the people and you sent Vani, the one on top
just because like she was a computer person and she liked to that kind of stuff.
At one point, one of the characters brought their daughter to work if she wants to make apps, and you chose Computer Science. You looked over the list for a while [which was not exactly true], do you remember why you chose computer science?
Just because you can make like apps on the computer and figure out how the stuff works
Can you tell me where you think apps in general come from?
mmmm
We talked about Minecraft earlier, if you think about in general, where do apps come from?
Probably from the people who make them, there's a whole bunch, people who made them, so, yeah
What kind of people do you imagine--can you tell me about these people?
Really smart and like really good at doing their job and working on computer stuff
Do you think that you could get a job making apps?
Ummmm. .... no. Probably not.
Why's that?
Just because I don't really like, think, about it a lot, and ... yeah.
OK. How do you think somebody would get that job?
Studying really hard probably
Do you have any questions about anything you saw, or any questions about the project or the study?
No.
\end{lstlisting}

\clearpage\section*{Subject 4}
\begin{lstlisting}
What is your favorite subject in school?
Um, English, like, writing. 
What do you like about it?
I like creative writing and being, like, imaginative
[offtopic]
What is your favorite computer app?
Probably like planner and calendar apps, like, organizing, so I know what's coming
Have you ever thought about the people who make those apps--the calendar and planning apps--have you thought about where they come from?
Kind of, because actually my teacher had us testing, she made an app which was a planner app such that she tested us with it. So after that I kind of did, but that was this year, but before that, not really.
[offtopic]
If you imagine the people that make apps, can you tell me what you think they look like?
Well, I'm pretty sure like people, they're like artists and stuff, like ... looks, I never really thought about that. [long pause] I picture like the people who _design_ it would be kind of like artists, kind of, messy kind of regular people, and the people who publish it would be kind of like, official, ... I don't know, yeah.
So these people who work on the apps, what do you think that they studied in college?
um...
You already mentioned there's kind of an art and design side, and a business side.
Yeah, yeah, um, art and design, and then, like, technology, if there's, like, a video game
And if you think about those same people, who make these apps, what subject do you think they liked in school?
Art... or technology
[introduce game]

"[on first page of intro]
OK, what do I, .. do I click?"
[affirm]

"[on second page of intro]
[chuckle]"

"[on third page of intro, gesturing to word ""system""]
What's this, like, your computer system?"
Yeah, the Social Jam product that they're making
[affirm]

[game start]

So do I, I'm not really sure what to do
[explain to click on alerts]

[explain task changing and clicking on faces]
Ok

Should I click on one of these here?
Sure, just before you do this event, keep in mind that you can click on these to change their tasks, too.
OK

So, are these other ones you can...
You can click on any of them
Oh

I'm not quite sure what [inaudible]
The indicator here tells you what their main job is
right
Vani is maintaining the system, keeping it secure, and Nancy is developing new features
Yeah... oh, OK, I got it

Can you tell me about how you're trying to make this decision?
I'm looking at what she studied, .. [inaudible], so what she wants to get um a job in and then like what the choices are.
[inaudible muttering to self about the options]

So how do we fix these [the exploits], is that what the alerts are?
By having people on the maintenance skill, they will be working on this.
OK
The maintenance task, I mean
Yeah

Do you want me to explain my reasoning?
Yeah yeah
OK, so, since the attack wasn't successful, and we need our people, [trails off]
Ok

[points out grammatic error in message]

[pointing to option to training employees in good passwords] Since we want to get better, so if the problem comes up again

[Indicating skill levels] So is this like how much they've done of these? Like, the numbers?
That's actually their skill at that, so how good they are at the task
Oh OK
[inaudible]

Do you know what a press release is?
No
[explains it]

Whats the line?
That's your goal.
Oh!

[clicks through the end screen]
That was the end of the game.
Oh!
It was the success screen, since you were able to  meet your goals.

[mute music]

Did anything in the game surprise you?
Um, it surprised me that I was able to get to the goal. 
Can you tell me why that surprised you?
I was expecting it to be um ... more ... i don't know [inaudible] I expected I'd have to teach more, the other people on there, um, more about [inaudible] because there's only like two or three, so that surprised me
Anything else surprise you?
[negative]
Can you tell me about what kinds of decisions you made during the game, the kinds of decisions you made and you how made them?
um, I... had to... decide whether to just let it go because it wasn't like a big deal or smaller than other problems or to um... ah ... there was some where you had to pick which of the workers would fit best with each thing, and ... 
I noticed when you had those kind of choices, you went in and looked at their bios. Do you remember what you were looking for, or how that helped you make a decision---or maybe didn't help you make a decision?
It helped me make a decision by telling me what they majored in, uh, what they were specialized in, and then just like how their personality is, like, the first one, I don't remember
Janine
Yeah, she how she was more at communicating [inaudible] other people, making them feel good, then other ones like the second one would more good at like designing and creative kind of things, and the last one was better at like fixing [inaudible] technology problems
mm hmm
So that just let me know, like if um they needed somebody to, like one of them with the hacker, they needed somebody to fix it, so I sent the last one because she was good at fixing that kind of thing. So, it just let me know which one would be best for each job.
Do you think you could get a job making apps?
Um, depending on what kind of app, if it was like a game, I probably wouldn't be very good at that, but if it was um..
Why do you say that?
I'm not usually on technology, um, usually like that, but like if there was like a coloring app or an artistic one or like a planner like I said at the beginning, I would probably be able to do that more designing. I would be better at that than [trails off]
OK, so when we imagine people who make apps, how do you imagine those people got their jobs?
Um, probably their teachers saw their ability to do something like that in an earlier stage, so then they were able to get the classes in like high school and through college to be able to get better at that
So, if a teacher saw that somebody had potential and guided them toward a path that ended up with them getting a job in a company like Social Jam, here in the game, what can you imagine that they would study in college or high school, what kinds of courses would they take? What kinds of work would they do?
I think I did, yeah, the art, and technology
Do you have any questions about anything that you saw, or about the study?
The music, it was good to have a background thing, but like it was a little bit distracting because I was trying to, like, read them. I'm not very good at, like, I always have to go back when I read anyways, but um, and...
There was a way to mute it, too, or we could have turned it down. I'm sorry. I didn't notice it was causing a problem for you. I have heard the song so many times that I have tuned it out.
... The bios were helpful to make decisions, like that, and yeah.
\end{lstlisting}

\clearpage\section*{Subject 5}
\begin{lstlisting}
What is your favorite subject when you're doing schoolwork
Maybe math.
Why's that
Just because it's fun to solve problems, Because I'm a problem solver, and because I like to know why it works, not just how it works.
What is your favorite app?
I don't have one because I don't really use iPhones or anything like that
What about laptop or desktop computers, do you do any work on those?
Yes. Probably Windows, PC.
What's your favorite kind of thing to do on Windows--what kind of programs do you run on it?
I don't know... Age of Empires is a game that is  a CD, stick it in the computer, but most of that I use for like schoolwork and stuff
You do some school work on the computer?
Yeah
Do you ever think about the people who make those programs
Sometimes
What do you think they look like?
I don't think about what they look like, I just think about how hard it would have been, and how long it would have taken, and how many people it would have taken.
How many people do you think it takes?
ha, a lot, over a long period of time
Well, what makes you think that---what makes you think it would be very hard?
Well, at least for Age of Empires, because it is, there is a bunch of like a whole lot of editing and like the, you know what the game Age of Empires is?
Yes
Well, just cause it's like, there's a lot of different things, and you have to make all the commands do the things, and stuff.
Do you have any experience with programming or making your own programs---coding?
Yeah, somewhat.
What kind of stuff have you done?
Well, I've done the Google's Code.org, and I've done some of GameSalad, some stuff on GameSalad Creator.
Great. My son uses one called GameMaker which is a nice program too
Is it free?
"Yeah, there's a free one, and uh we're going to start using another one soon called Construct 2 which a student had recommended to me, so I'm going to check it out.
So, when you think about these people who make these apps, whether they're making things to help people do schoolwork, or... instead of 'apps' let me just say 'programs', people are making Age of Empires or people are making programs you use to help you do your schoolwork, what do you think they studied in school?"
pffffooo. Lots of stuff. Animation, first off for a lot of them, and programming, and stuff. Like, and they would math and stuff.
OK, anything else?
uh, well, all the basics, well everybody does, for the most part.
Well the beauty of being homeschooled is that you get to define what the basics are, so what do you think the basics are?
That's true.
That's why we homeschool too---we can define the basics for ourselves.
English, uh English, Math, Science, uh... well, we say Bible too, and then ... I feel like i'm missing one. I forget what it is.
OK. Want to try the game out?
[explain think aloud, and starting the game]

Do I click to...?
mm hmm

[end of intro]

I have to make the.. do stuff?
mm hmm
... so, what do I do?
You can start by clicking down here, on the alerts button

[explain how to change tasks and clicking faces to learn about workers]
[he goes through and reads about each worker]

"[has expanded bottom worker, so the main interaction area is not visible]
So then, I don't exactly understand what I'm supposed to do"
If you click on Melissa again, the other screen will come up, and she [Frieda] will tell you
OK

Do you know what a press release is?
mmmm no
[explains it]

[After sending one worker to deal with a different task, probably security conference]
Can you tell me why you switched them both to maintenance just now?
Ah, because the one is gone, so [inaudible]
OK

"[Looking at daughter advice screen]
Is it just one of them?"
You can only choose one, right. Can you tell me about how you're making your decision?
Um, because I think it would either be one of these two [CS and Graphic Design], because this [CS] would be more of programming and how to work the, and how to work the computers, and this [graphic design] would be more like the designy parts of it, so I think it could be either of these. Ah... let's just say... [clicks Computer Science]

"[He is looking over the event where 'report to police' is an option]
Can you tell me about the decision you're trying to make?"
Um, no, I don't really know. Um... don't really know, it's either one of these two [ignore/report], I think, but I don't know. ... Hmmm... I don't really know, so I think I need to say report to the police. 

So, how do I make it go down?
The chart?
The, yeah
The chart actually is showing your, if you click on status you can take a look at it while the game's paused, this is showing how many users you have, so you actually want that to go up
OHH!
You want users up and you want exploits to be low.
I see. I see. I was, I thought that this [progress toward goal] was the one you had to keep below 15%.
Ah, no, the estimated exposure, if you click on status here [he does], it's the estimated exposure you want less than 15%. This one [progress toward goal] you want to go up.
The.... exposure... I see, OK, OK. I'm like, [adopts panicked voice] "Oh oh how do I make that go down?! Cause I just lost!" OK.

[talks to himself, inaudibly. he does this occasionally, often when reading on-screen text]

"[Facing the password event]
[laughs] Encrypt it [laughs] That's not one of the options. [laughs] Would this [train] be training to encrypt it, basically?"
Well, that's an interesting question. Encrypting is something you do on the server-side, to protect your data, but it's not something that guy... guy in advertising, something that they can do directly.
Right.
So, yeah, so those are the two choices that they give you.

Can I look them up and then get back to this?
What's that?
Can I look at them and then get back to this?
Yes, yes.
OK

What's PPH, or , PHP? [Vani's interest]
PHP is a programming language for making Web sites.
OK.

"[ Looking at end game screen ]
Are exploits and features good or bad, what are these?"
Features are good and produce users, but the exploits cost you users.
Mm kay.

[game ends]

When you played the game, did anything surprise you?
Well, it first, uh, where it's uh of the thing it says you just got attacked, whatever, and then, after that it wasn't really surprising, it was just like, OK, I got it, whatever
So you're saying---I want to make sure I understand---the first time you were attacked by hackers, you were surprised, and then after that you kind of expected them to attack some more.
Yeah, I tried to get ready for it, ... and then, uh, it's I actually kind of liked it, I liked the way you have to try and run your company, I just kind of like that type of thing. I don't know about other kids, but [trails off]
Sure. Anything else surprise you in the game?
... I don't think so
Can you tell me about what kinds of decisions you had to make while you played?
Whether or not to have somebody work on, like, whether or not to make things to create, uh, send people to, uh, to get trained so that they're more effective, then whether or not to have some try to find the hackers, and whether or not to, uh, whether I thought I could handle it by myself or talk to the police or whatever
Do you remember when you were prompted, you had the option to send someone to a security conference, you sent Nancy. Do you remember why you sent Nancy?
Because she was the one that was, I don't remember what it was, because of when I was reading about the people
In their descriptions?
In their descriptions. 
Sure. So do you think that you could get a job making apps?
Yes.
Why do you say that?
Um, well, if I'm, ... you mean me personally?
Yeah
Um, maybe, yes probably because I like to, uh, I can, ah, work on like programming stuff [inaudible] programming.
How does someone get a job making apps? How do you imagine they end up having that job?
ummm, through an interview [chuckle]
What do you think leads to the interview---trace backwards from there
Well, somebody would hear about you, and uh, you probably need to get a degree, or be an entrepreneur and start your own app-making business
Mmhmm, "entrepreneur" is a big word for a twelve year old---that's a good word.
Antidisestablishmentarianism.
[both laugh] That's a good word too! Do you know anybody that's started their own business?
Well, me, my brother---well, me and then my older brother, the one who's not here, uh, [inaudible] friends, yeah.
And you started your own business
Yeah, it's, we, uh, got a little machine that takes pictures of old slides so that I can, people will send in a box or whatever of slides and I'll take pictures of them and convert them--or I'll take pictures of them and convert them to digital, and I'll give them a CD with them on it for 25 cents a slide
Great, that's cool
Do you have any questions for me or any questions about things that you saw?
Mmm not really, um, the one thing is, just at the very, at first, if I were playing the game, if I didn't have anybody here, at first, I wouldn't know what to do, because I saw the thing flashing, but I didn't know you were supposed to click on it.
Sure. There's some places where it's a little rough around the edges, it's not quite ready for public release, and that's one of the reasons we're bringing people to play it early.
\end{lstlisting}

\clearpage\section*{Subject 6}
\begin{lstlisting}
What's your favorite subject in school
I'd say lunch but that probably doesn't count
I'll take lunch!
Sure.
Or you can list more than one, it doesn't have to be a single one.
I don't, I'm not a big fan of school
That's allright
I mean, I know I have to do it, but it's not very fun
Sure
What's your favorite app?
Favorite app...
or favorite computer program generally.
Uh... I don't know, I like Minecraft. It's fun.
What do you like about Minecraft.
I don't know, it's just fun, I'm creative so I like building stuff. So... I don't know.
So, if you think about a program like Minecraft---started by one guy, you probably know, and then now it's being worked on by a very large team of people
right
have you ever imagined those people that do that work, the people that work on minecraft?
Mm, well not necessarily Minecraft specifically, but I've thought about that kind of thing
What do you think they look like, the people that work on that?
I don't know. Probably, I don't know. Probably most of them are not American, probably a lot of them are from Japan, it's a Japanese company, isn't it? And then, I don't know. Probably just mostly that, just a bunch of geeky guys that sit at their computers and program Minecraft all day.
OK, and do you think they went to college and what do you think they studied?
It probably depends on what ones, and they probably studied, like, animation or, I don't know, something like that
And so these same people that you're imagining, what do you think their favorite school subjects would have been?
Probably art or math.
OK.

[introduce the game]

[finish intro]

[ explain task changing, tap on faces to learn more about them ]

[ mumbles to himself while making some choices, but nothing clear ]

Can you tell me about how you're making this decision?
I don't really know, I'll just kind of guess, I could report it to the police, or I could task someone to try to figure it out. I'll probably just task somebody to figure it out for now because, I don't really feel like getting mixed up with the police, so I'll task [someone] to do it.

OK, so I should not actually have had him do it.
Why do you say that?
Because he got a degree in electrical engineering and he got a degree in computer science, so... [trails off] [tries changing the task]
Yeah, once he's on a special task he cannot be taken off of it.
I can't... OK

OK, tell Jerry to do it.

Hm, interesting

OK, so I have to figure out... how to fix the problem. Somehow.

Interesting, it's raising faster.

Oops, I should have sent...
What were you going to say?
I should have sent this guy.

Ok, so, that's the end?
That's the end
OK, cool.

Did anything in the game surprise you?
uh, yeah, a little bit. It kind of surprised me the random losing the customers because of the stuff popping up. So, then, I was kind of, like I didn't really know what I was doing, so I was trying to figure it out, and ... so then, it was like, I hadn't researched, I hadn't looked up who, I hadn't learned anything about my employees at the beginning, so I accidentally told the wrong guy to do something, uh
And that was, you said before, you thought it was the wrong guy because he had the electrical engineering degree instead of computer science?
I mean, he still got it done, but yeah, and then, yeah, I guess. So that was pretty much it.
When you had options to send people to the security conference, you sent Jerry, the guy at the bottom. Do you remember why you sent Jerry?
Yeah, I sent him just because he was uh he was one of the uh um computer science people. I figured he might be more helpful to, uh, computer science guys at a security conference.
Sure. And I noticed too when James from accounting brought his daughter to work, you looked over the options and you chose Computer Science there too. Do you remember why you chose that?
Uh, just because it seems like that's what, it seems like making an app would be more computers and that stuff than anything else.
OK, Tell me about where apps come from.
Anybody that makes them.
Can you say a little bit more?
So if somebody decides they want to make an app and they know, they sit down and do it. If they don't know how, and they still want to do it, then depending on how complex the app is they either go to try to figure it out at some random source, or if it's a bigger, more complex project, they go take classes for it... and then they make the app.
Do you think that you could get a job making apps?
Probably, if I wanted to.
Why do you say that?
cause if I learned how to do stuff and then, I don't know, applied for jobs, especially if I got a degree, because people would be more likely to want someone with a degree in app development, or something along those lines
Do you have any questions about anything that you saw or any questions for me about the study.
[negative]
\end{lstlisting}

\clearpage\section*{Subject 7}
\begin{lstlisting}
What's your favorite subject in school?
I like English and science
What do you like about those?
You get to be more creative than in other classes
What's your favorite app?
[perplexed]
Do you have a favorite app? Do you use any computer programs that you like?
Mmm, not really, I'm ... I've been trying to learn how to like develop games, but [trails off]
What kind of work have you done along those lines? Or what have you been reading?
I've been, I've read homestuck, and I've been inspired to make a webcomic that I don't have a name for yet, but I'm writing the script for it.
[offtopic homestuck]
Can you try to imagine the people who write computer software, who write apps.  Can you imagine what they might look like?
[perplexed]
Or could you describe how you imagine they would look?
I imagine them to be like 17 to sort of early 20s in age, but appearancewise, I can't really [trails off]
Can you imagine what subjects they would like in school?
Probably English and Computers
Why do you say those?
Because in English, you get to learn how to develop characters and plot and stuff, and Computers, obviously, you learn who to deal with computers.

[introduce game]
[slide two, chuckles at the protect from hackers]
"You have three employees"--that's not a lot! [chuckle]

[ explain task change and tap to learn about employees ]

"[DDOS event]
Angry customer?!"

Can you tell me why you chose Jerry?
He seemed like with all, well, hold on, [opens his profile] "He worked for a local company in high school and has worked as a software developer ever since" so it seems like he would have the most experience and .. most, like .. willingness I guess.

"[looking at exploits screen]
[hushed shout]Fix it!"

"[Script kiddie event]
Just a script kiddie, so [trails off]"

[Spends some time looking at exploits list]

[angrily, hushed] Gosh, so many alerts.

[end] Yay!

Did anything surprise you while playing the game?
I think maybe the amount of exploits and all the choices, sort of, that you had to make.
And can you tell me why that was surprising to you?
That was, I mean, I know that starting up anything there's going to be some problems, but I didn't expect that many that quickly.
When you had the choice to have one of your employees investigate the injection attack--it was kind of toward the end of the game--you chose Jerry pretty quickly. Do you remember why you chose Jerry?
Yeah, well he seemed the most well-trained out of all of them, and the girl--I forget her name--she seemed more like someone who should be with, like, social, and not developer, but Jerry seemed like the better person for the job of that kind of thing.
How would you describe the decisions that you had to make in the game?
You really had to think of what sort of impact it would have later in the game when you would choose them, but you also had to make them kind of quickly, so... [trails off]
So how did you make those decisions? You mentioned thinking about future impact.
Thinking about how many people you could have lost, or how many people you can gain, from certain choices you take, so [trails off]
Do you think that you could get a job making apps, like at this company?
Yes.
Why do you say that?
... Like, it feels like, uh ... the way the company was going it was going on like strong and reliable and being able to tell users what's going on,  so... [trails off]
I don't disagree that you could work there, I just wondered what you thought about it. So, how does somebody get a job making apps?
They would have to learn how to code and script and stuff, and be able to like produce a game that people would play with it, and be able to develop a plot and a story, or whatever you needed for that game.... while still being semi-entertaining. 
Do you have any questions about anything that you saw, or any questions about the study that we're doing?
[negative]
\end{lstlisting}

\clearpage\section*{Subject 8}
\begin{lstlisting}
What is your favorite subject in school?
Math.
And what is your favorite app?
Tererria.
Oh thats cool! Great game. So who are the people who made Tererria?
Like...they have huge computers. Spend a lot of time working. Oh, and there have to be some super talented artists. The people probably went to college for like engineering.
You said engineering. What do you think their favorite subjects were when they were in school?
Uh...Probably science or math. Mainly that.
[introduce game]
[game play notes] He actually read all of the text. There were a lot of events from the get go so time took a long time to actually move (4 events in day 1).
[He took time to decide whether to do nothing or report to the police.]

How was it?
It was fun and different.
Was there anything that surprised you?
Yeah when the guy asked...I mean the guy's daughter asked about what she should study, I was surprised there wasn't the option to choose more than one thing.
Oh and what all would you have chosen?
Computer Science like I did. And engineering and math.
Ok that interesting. What kinds of decisions did you have to make?
Uhm...
Rather what were the choice you made and how did you make them?
I had to choose, like, the skill to have people on, and I used their degree. What they were best at. The other things I thought about what was good or what would get the most...points...Users. Yeah users.
So again, where do app come from?
Code and computer programs. And I guess the people who write the codes. Computer scienctists and engineers.
Do you think you could get a job making apps?
Yeah I do.
Why?
Well I've coded at my school at the beginning of the year and I was pretty good at it.
Oh so you've coded?
Well not really coded. It was like where we had to put things together...
So block coding?
Yeah, that!
What grade are you in?
Sixth.
Thats very cool. So how does someone get a job making software?
Probably get a degree. Spend four years in college. And arists...would have to practice a lot. Spend time getting better. And become tech savvy.
Ok. So lastly, do you have any questions for me?
No.
All right, thank you so much!
\end{lstlisting}

\clearpage\section*{Subject 9}
\begin{lstlisting}
Ok so thats recording. So I'm just going to ask you a few questions to start out with. Whats your favorite subject in school?
Uh...academic?
Yes.
I'd say english.
Ok and then just broad favorite subject in general?
Art. Or band. Either one.
Ok and thats visual art? So like drawing and painting.
Yes.
Whats your favorite type of app?
Drawing.
So like you have...?
I like sketchbook pro. Its my favorite.
And then who are the type of people who make sketchbook pro?
Adobe I think.
And if you think about the actual people who made that?
The people? *pause* I don't know...
Ok yeah, thats a valid answer.
Yeah I've got no idea on that one.
So no idea like maybe what they would look like or what they studied in college if they went to college?
They probably went to college if they are making such a large app. And working for adobe. They probably al crazy mad scientist. Like...Muahaha. We are gonna make this crazy app so people can draw.
Ok and then what...when those people were in school, what do you think their favorite subjects were?
Probably math and computer sciences...maybe art. Yeah.
Ok.
[Introduces Game]
[First Play Through]
The music is groovy.
Do you have what it takes? Yessss
Idk, I really like nancy. Soo I think I'm going to go with...oh Oh, what do? Oh, ok.
So she has a bachelors of english...oh my mom when to Georgetown! Nancy has a popular podcast about being a woman in technology. Oh thats cool!
Bruce has a... [reads text]
Ok lets see what Janine has for us. [Reads Text] Ok! I think I'm gonna go with Janine, because she sounds really cool.
So all three of them are your employees.
Yeeeahh I like them. They're cool.
Ok so status...features...oh selfie of the week!
[reads instructions]
[does not reassign anyone]
Ok. So. So we still don't have a selfie of the week!
Oh no, 55 users left....no! [notify users option response]
Oh. Oh no. You were attacked... [reads text] Oh *frantic* I'm going with! Ok. ok *frantic* uhm....especially talented at meting with customers...has a popular podcasts... and he...ok. I think I'm gonna go with Janine, because she is my favorite. So...or. Or report to the police. Because thats probably something I should do. 
The police found the hackers! Free publicity! Wow thanks hacker guy.
Our users are really going up. This is going great.
Oh oh! I'm really liking this. I have a lot of users now. They're really liking my game. You don't even know how happy this is making me. 
Oh! Oh no. An angry customer... [text] Just gonna wait this out. NO? After two hours...no! So many users, 437 have left because...jeez. sorry. 
I would say computer science. 
Oh one of your workers can go [text]. Ok I'm gonna send bruce because he sounds like he spends too much time watching history documentaries.
James from accounting called [text] Oh wow. That makes me feel so good about myself!
OK I'm gonna send Nancy because Janine is my fav and we uh need her handy.
Oh Bruce came back! With greatly increased maintance skill. Oh wow that is greatly increased.
Uhm...lets train the staff because why not.
Because you could fit..uh...what...no....ok. Its ok, at least we fixed it.
Two hundred...left..no stop leaving me! Thankfully we took time train our staff. Good! Good it wont happen again. I need my users.
Oh mannnn oh I really didn't get enough users.
Ok one second. *aside* So I don't need to clear the logs? Just restart, correct?
Yes, we will just append the logs.
Cool!
Ok so I'm gonna show you something at the beginning that will make the game easier for you.
Ok! I like Nancy, she is sassy. All of them are sassy.

[Start game 2]
Oh! New people. I like this guy. (Ivar)
Alright so if you go down here. And click this. Here Frieda says you can reassign any of the to development? So click here, yeah, and you can put them on development. So last time...
They were all on maintanance.
Correct, so you weren't actually making any new features.
Wooah this is crazy. I like Ivar, he looks kinda creepy and I love it. I like his creepy long hair with a bald head.
So you see how the features...yeah the progress is at 42%? Its going up because its actually being made.
Oh wow! Wow look at that...tumblr...wow.
One of your workers can go to a two day conference...no I'm not gonna send Ivar because I don't want him to creep out all the...I'm gonna send Vani.
[Feature Completed] Oh wow, so many people. They love my game! Or...whatever we are doing. I think its a game. Yeah its a game!
Oh report this to the police. The police will help.
Vani, yes, my fav! Great maintanance, wow. I'm changing you to development.
And I reached my goal!
Hackers stole some users data. Ohhh Ignore it! Phew avoided that. -- An independent securiy...wha---No! Maaaan.
Its ok. Its ok.
An angry customer....I'm just gonna wait this one out. Nooooo come back! Ok. Ok.
This is intense, honestly.
One of our engineers....ok. Estaban. I'm gonna have you doing that.
Ok, train this staff. And now everyone left because they're like, "Wow. They're dumb. Don't have good taste in game."
Aghh almost there. Just did a lot of bad exposure things
[end game]

Ok so I just have a couple more questions. Did anyting surprise?
The first time I didn't get the development thing. That really changed it.
Anything about the story or the characters?
I really liked it actually. Nothing that like surprised me that completely caught my attention. It all seemed pretty natural, nothing seemed out of place.
Thank you. What types of decisions did you have to make in this? And how did you, uhm, decide. How did you make those decisions?
Uhm when the hackers were on me, I had to decide what I'd let the public know, and what I didn't want to let the public know. 
Mmm hmm.
So I thought that was kinda cool. I really liked how I had to choose which one went and did the jobs.
So having gone through this a couple times, how did the decisions pan out from when you let the public know to when you didn't? Did one end up being better?
See, last time I was like "oh I'll let them know, and everyone got mad." So I thought I'd just keep it a secret. But then when I was exposed, everyone was like, "why didn't you tell use that." They all got really mad and left me, which I think is what made my exposure go up. But last the didn't like that. This time they don't like it more. 
Yeah. Ok I'm gonna stop this music from playing. So where do apps come from/
Uhm, computer developers and companies. A lot of them come from apple itself. Because, ya know, the apps that are already on then. Yeah, coding obviously. You can't really make a computer program witout coding.
Do you think you could get a job making an app?
I think I could, yeah! I seem to be pretty good at coding in the coding classes we have at school. We have on called B.I.T. I don't remember exactly what it stands for, but you do a lot of coding in that. 
Thats cool.
I think I totally could.
How does someone get a job making apps?
I assume that after they get a degree in computer science, or I'm sure there a school that specializes in like app creation. You could go to like adobe, like the ones that made the sketchbook app, uhm, if you're looking for more artistic or graphic designee stuff. Or you could go to apple. Game companies, develop applications there. 
Is there software that exists outside of game companies?
Yeah like...uhm...scheduling and planning. And like...thats the only one I can think of because I don't use them much. OH music! How could I forget that. There are a lot of music applications.
Ok so the last thing, do you have any questions for me?
No.
Cool. Well then--
Let me throw in a question from left field here. What is it that made you like Janine so much?
Uhm. She was idk. From her description she seemed really on it. She seemed really cool. I liked her a lot.
And at the very beginning, the first thing you said was that you liked, uhm--
Nancy?
Yes, Nancy. Was there a reason? That was before you had openend up and read about her.
She looked really sassy. And she looked like me before I went insane. Cuz I used to have long blond hair and that was normal. It was gross. But she was really sassy so I liked her. Sassy like not really sassy but like, cool. Like she looked hip with the kids.
Yeah! Ok, I think thats all we have. Thank you!
\end{lstlisting}

\clearpage\section*{Subject 10}
\begin{lstlisting}
Ok So your favorite subject is?
History.
Whats your favorite app?
Facebook.
Ok. Any others?
Mostly social media. Facebook, Instagram...I have this called wattpad I used all the time.
Where do the apps from?
From the app store. People make them and sell them on the app store.
What do you think those people are like?
Like they know how to program stuff. And...they know how to use the technology to build the apps to put them on the appstore for people to use.
If you had to describe the way that they look?
The way that they look? Uhm...smart? Or...I don't know. They look smart or techy I guess.
Do you think they went to college?
Most likely. Probably.
What do you think they studied?
Uhm, I'm not sure.
Thats ok! Valid answer. 
[Introduces game]
>>5:30<<
[end game]
So. What surprised you in that?
Well there was like, a lot of different...ok. It was different from any video game I've played. I had to think about what I was doing before I just clicked anything. I had to figure out what this button meant before I started just doing something. I had to kind of learn as I went.
What types of decisions did you have to make?
Well I had to decide whether I just reset the password or I would train everyone on how to decide on passwords so that they wouldn't get broken into again. And then I had to decide which person was going to do the maintanance part and which person did the development part. And that probably affected...something.
Yeah when you...did you see...you saw the features page and how the percentage was going up?
Yeah.
The people working on that were the feature people, and the people on maintance were on the exploits page, the things taking away users, the people on maintance were getting rid of those. So how did you make those decisions. 
Well I've always thought about not having weak passwords, so for the password one, I always try to make them as strong as possible so no could find them out. And for the other ones I just experimented for the first minute to figure out which was the better one.
Specifically about the one where it said someone had attacked your company and you had the option to retaliate, why did you decide to realiate?
I guess I'm an agressive person. I was kind of..
Was it good?
No, I probably shouldn't have. I probably should have let the police handle it instead of taking it into my hands.
Maybe. Good to know. Where do apps come from?
People like you who make them and develop them. Take a lot of time to make them.
What do those people look like?
Like you? Like normal people.
Just like me?
Like anybody!
What did you notice about the characters?
I noticed some of them had degrees from Ball State. And then they didn't look like smart people, they looked like normal people.
Do you think you could get a job making apps?
Possibly, if I thought about it and was patient with myself, then probably.
Why?
I can because I don't have to be super smart. you still have to have some knowledge about how to program and everything, but I think once I learned how to do that I would have a little bit more...uhm...chance of being more successful and making an app or a game.
So after having practiced?
Yeah!
Alright so last thing, any questions?
Nope! It was fun.
\end{lstlisting}

\clearpage\section*{Subject 11}
\begin{lstlisting}
All right, I'm just gonna start with a few questions. First off, what is your favorite subject in school?
Hmm...probably Gym. Does it have to be like an englihs or math?
No, thats good. But if you had to pick an academic one?
Probably science.
Alright, whats your favorite app? Which one do you use the most.
I don't actually have my own personal device like that, but we have a family one, so...
It could be, like, whats your favorite game? If you play video games?
When I do play, probably Clash of Clans.
Alright. Where do apps come from? Like apps or games.
The app store? Or do you mean before that?
Uhm. Yes, but, if you think of the people who made them.
Well I do read who made them, if its like catch app or someone like that.
If you think about the individual people who work on them, what would you say they look like?
Like, physical appearian--appea--
Mhhmm, appearance.
Uhm.
And so you know, I don't know is a valid answer throughout this whole process.
Yeah, I don't know.
Do you think they went to college?
I would assume so to learn how to make it. But I have some friends who can program stuff who haven't been to college, so.
If they did go to college, what do you think they majored in?
Probably some sort of technology...computer science maybe.
Alright, thats all I have beforehand.
[introduces game]
Whats a press release?
Its--like contacting a news station, or I guess today using social media. Saying what happened, letting the public know what is going on. 
Ah ok.
[end game]
Alright, do you want a second go on that? You figured out half way through that you can change people's task, but that really affects the game. 
Yeah, I'd like to give it a try.
Alright lets start you off. So I'm not sure what happened last time, but right here are the instructions, so you can read through those.
[Start Game]
[End game]
So did anything surprise you?
Well, a lot of the video games I play aren't very realistic. This one is pretty down to earth. In like the things that happen. And up to date.
So, what types of decisions did you have to make.
Well you had to work with your customers and learn how to make them happy while also...uhm...increase your benefit and keeping away people who wanted to steal that.
How did you make those decisions?
Well I knew some of them would hurt me no matter what I did so like the best possible answers. And some of them, I really didn't know what to do, so I just guessed.
Do you remember any in particular you just guessed about?
Uhm...The one, the first problem. Where you could tell the police about it or send someone after it.
The first time you played through you chose to have someone retaliate and the second time, you reported it to the police.
Yeah, neither of those worked, I probably should have just done nothing. --- I don't really remember what happened with the police, actually.
You celebrated because they found who it was and you got users because of the publicity.
Oh! I thought that was a different...ok nevermind! I thought it was a different one. Yeah that one worked, so I went with that.
Same question from before, where do programs and apps come from?
Uhm, yeah people, groups of, like, you buy them from the app store, but people will make them on their own. Or different organisations. 
Do you think you could get a job making apps?
I could but thats...thats not really...I don't know a whole lot about it. I could if I took that career path.
What is that career path? How do you get to that?
There was computer science on there and that worked both times I tried that, so probably something like that. Thers probably a lot of mathematics and, you know, how you have to know how games work or what other people did to make the games.
Yeah so how do people in general gets jobs in making software? Like do you think they went to college and studied and?
My uncle has a job in computer software but he didn't graduate from college. So I don't think you have to.
So there are multiple ways?
Yes. 
Alright I think thats all...one more. What did you think about the characters? You read their bios several times. 
I thought that was probably a pretty balanced team, because you had someone who, like, majored in three of the important categories, such as computer engineering. I can't remember what the one in the middle was, and then there was another one who was good at working with people and finding out what they want. Thats why I chose her to do the--oh no, I chose the engineering one to teach people how to do the password thing because he knew how to make a safe password. And then I would have chosen her to find out what people liked or stuff like that.
Alright thats all I have. Do you have any questions for me?
Not really.
\end{lstlisting}

\clearpage\section*{Subject 12}
\begin{lstlisting}
I'm going to start off by just asking you a few questions. Whats your favorite subject?
Writing.
Alright. What is your favorite app or game? On technology.
I don't know.
Where do apps and games come from?
A creator.
So a person makes it. What...what do you think that person looks like? If you had to imagine the person who makes these things, what would you say they look like?
I don't know, I've nevers seen them.
Ok. Do you think that they went to college?
Probably.
If they did, what do you think they studied?
Computer technology.
[introduces game]
Is it just gonna move on, or?
Just click on the screen.
*time passes*
So every time this flashes, time stops and it has someting for you.
*time passes*
What is...?
Ok, if you click here, you'll see...right there, 28% on development? If you put someone's task to that, they work on completing it. That percentage will go up. Once you take care of this event, that'll start goin up again.
*time passes*
What do these do?
Uhm...so now time has stopped since there is something here. Click on this one. These are are things that take users, but since you have him on maintance, this one is already fixed. Now click here. These are things you are making. She is working on those and once you have that you'll start making more users.
[end game]
Did anything surprise you?
Not really.
Ok. What types of choices did you have to make?
Uhm. Like whether or not to train staff and have them not working for a time.
Mmhmm. And how did you make that decision? How did you choose whether or not to do that and who to send?
I don't know...
Ok! Thats fair. After having played, same question, where do apps come from?
The creator.
Do you think you could...when you're grown up, do you think you could get a job making software--making apps?
No.
Why not?
Because its something that interests me so its not something I could work very hard at.
Ok. What do you find interesting?
*confused silence*
Another way--what do you want to be when you grow up?
Really, only God knows what I'll be and I, I don't know.
Ok! Fair. In general, how do you think someone gets a job making apps?
You would go to college and study and...get hired or start a business.
Mmhmm. Study what?
Computer science.
Alright. Those are all the questions I have for you. Do you have any for me?
*Shakes head*
\end{lstlisting}

\clearpage\section*{Subject 13}
\begin{lstlisting}
To start, I'm going to ask you a few questions. What is your favorite subject?
Uhm...writing.
Alright. What is your favorite app?
What?
What is your favorite app? App, game...so like on technology, do you have a favorite? Its ok if you don't.
You mean like a game or anything online?
It can be anything really. If your favorite happens to be the dictionary app, thats alright too!
I don't really have one.
Ok. Thats fair. Where do apps come from? 
Well, people have to make em.
Mmhmm. Describe those people.
Uhm, people that know computers. Computer programs.
If you had to describe what they look like? If you imagine that person in your head and you were describing what they physically look like. If you were gonna write about them?
Uhm...geeky...
Geeky?
Yeah.
Alright, do you think they went to college?
Yes.
What do you think they studied?
Computer science. 
[introduces game]
Feel free to ask questions, I'll explain anything you don't understand.
What do I do?
You just click.
What do I do?
Click on the button that is flashing.
Uhm...
So if you click on him here...yeah. So you can do that for any of them. And then if you click on these buttons--
Here?
Yes. There you can decide if you want them on developing or maintaining because thats what they are working on.
And now what do I do?
Click on him again and that'll make it smaller. Then click here.
The red?
Yeah, whenever that is flashing like that, it needs your attention.
[end game]
Ok, did anything surprise you?
Not really.
Ok! There are no wrong answers to these questions. What types of decisison did you make.
Uhm whether to train the staff.
Mmhmm. How did you make that decision?
Train em.
Yeah thats what you decided. Why did you decide that?
To....secure the area more. Company more.
Alright. Where do apps come from?
Where do apps come from?
Mmmhmm.
Well...well still people have to make em still.
Do you think you could get a job making apps?
No.
Why not.
Well I guess I could but...
But?
Well I'd have to take the right school stuff.
So you could but you would have to study the right things. What would you have to do, to study in order to?
Computer engineering and computer science.
Ok! Is that something that happens only in college or can you start learning beforehand?
You can beforehand.
How does someone get a job doing that.
You have to be experienced enough to...for someone to want to hire them.
Ok. How...after playing that game, how important is it to make sure that you're doing on computers is safe?
Very.
Ok! Thats all I have. Do you have any for me?
No.
Alright, then we are done! Thank you very much for participating.
\end{lstlisting}

\clearpage
\section{Gameplay Logs}
Below are the gameplay logs from each research subject, identified
with the subject's number.
When a session included more than one play through the game,
the logs are provided sequentially and annotated with a sequence number;
for example, \verb+1-2+ is subject~1's second play of the game.

\begin{verbatim}
1-1:
CEO is Nancy; workers are Janine, Vani, Abdullah, 
edu.bsu.cybersec.core.narrative.InsecurePasswordEvent: Train staff on secure passwords
edu.bsu.cybersec.core.narrative.ChildAdviceEvent: Graphic Design
edu.bsu.cybersec.core.narrative.DDOSEvent: Press Release
Set volume to zero
Set volume to zero
edu.bsu.cybersec.core.narrative.SecurityConferenceEvent: Janine
edu.bsu.cybersec.core.narrative.InputSanitizationEvent: Vani
edu.bsu.cybersec.core.narrative.DataStolenEvent: Notify our users
Final users: 2603; exposure: 0.0069495

1-2:
CEO is Melissa; workers are Esteban, Ivar, Bruce, 
Esteban task changed to Development
Esteban task changed to Maintenance
Esteban task changed to Development
Ivar task changed to Development
Ivar task changed to Development
Esteban task changed to Development
Expanded: Ivar
Collapsed: Ivar
Esteban task changed to Maintenance
Ivar task changed to Maintenance
Bruce task changed to Development
edu.bsu.cybersec.core.narrative.ChildAdviceEvent: Graphic Design
edu.bsu.cybersec.core.narrative.InputSanitizationEvent: Ivar
edu.bsu.cybersec.core.narrative.InsecurePasswordEvent: Train staff on secure passwords
Esteban task changed to Development
Esteban task changed to Maintenance
edu.bsu.cybersec.core.narrative.DDOSEvent: Just Wait
Bruce task changed to Maintenance
edu.bsu.cybersec.core.narrative.DataStolenEvent: Notify our users
Ivar task changed to Development
Esteban task changed to Development
edu.bsu.cybersec.core.narrative.SecurityConferenceEvent: Bruce
Final users: 13112; exposure: 0.1412592

2:
CEO is Abdullah; workers are Janine, Melissa, Esteban, 
Janine task changed to Development
edu.bsu.cybersec.core.narrative.DataStolenEvent: Notify our users
Esteban task changed to Development
Melissa task changed to Development
edu.bsu.cybersec.core.narrative.InputSanitizationEvent: Janine
edu.bsu.cybersec.core.narrative.InsecurePasswordEvent: Train staff on secure passwords
edu.bsu.cybersec.core.narrative.SecurityConferenceEvent: Esteban
edu.bsu.cybersec.core.narrative.ChildAdviceEvent: Computer Science
edu.bsu.cybersec.core.narrative.DDOSEvent: Press Release
Final users: 26191; exposure: 0.0674723

3:
CEO is Jerry; workers are Vani, Bruce, Nancy, 
Expanded: Nancy
Collapsed: Nancy
Expanded: Bruce
Expanded: Vani
Collapsed: Vani
Nancy task changed to Development
Bruce task changed to Development
Expanded: Vani
Collapsed: Vani
Expanded: Bruce
Collapsed: Bruce
Bruce task changed to Maintenance
Expanded: Nancy
Collapsed: Nancy
edu.bsu.cybersec.core.narrative.SecurityConferenceEvent: Nancy
Expanded: Vani
Vani task changed to Development
Vani task changed to Maintenance
Bruce task changed to Development
Collapsed: Vani
edu.bsu.cybersec.core.narrative.DDOSEvent: Press Release
edu.bsu.cybersec.core.narrative.InputSanitizationEvent: Vani
edu.bsu.cybersec.core.narrative.DataStolenEvent: Notify our users
Vani task changed to Development
Bruce task changed to Maintenance
Nancy task changed to Development
edu.bsu.cybersec.core.narrative.ChildAdviceEvent: Computer Science
edu.bsu.cybersec.core.narrative.InsecurePasswordEvent: Train staff on secure passwords
Final users: 6725; exposure: 0.092293896

4:
CEO is Esteban; workers are Janine, Nancy, Vani, 
Expanded: Janine
Collapsed: Janine
Expanded: Nancy
Collapsed: Nancy
Expanded: Vani
Collapsed: Vani
Nancy task changed to Development
Expanded: Nancy
Collapsed: Nancy
edu.bsu.cybersec.core.narrative.DataStolenEvent: Notify our users
Expanded: Janine
Expanded: Nancy
Expanded: Vani
Collapsed: Vani
edu.bsu.cybersec.core.narrative.SecurityConferenceEvent: Vani
edu.bsu.cybersec.core.narrative.ChildAdviceEvent: Graphic Design
Expanded: Janine
Expanded: Nancy
Collapsed: Nancy
edu.bsu.cybersec.core.narrative.InsecurePasswordEvent: Train staff on secure passwords
edu.bsu.cybersec.core.narrative.DDOSEvent: Press Release
Expanded: Janine
Expanded: Nancy
Expanded: Vani
Collapsed: Vani
edu.bsu.cybersec.core.narrative.InputSanitizationEvent: Nancy
Final users: 39247; exposure: 0.05003875

5:
CEO is Ivar; workers are Nancy, Vani, Melissa, 
Nancy task changed to Development
Expanded: Nancy
Expanded: Vani
Expanded: Melissa
Collapsed: Melissa
edu.bsu.cybersec.core.narrative.DDOSEvent: Press Release
Melissa task changed to Development
edu.bsu.cybersec.core.narrative.SecurityConferenceEvent: Nancy
Melissa task changed to Maintenance
edu.bsu.cybersec.core.narrative.DataStolenEvent: Notify our users
edu.bsu.cybersec.core.narrative.ChildAdviceEvent: Computer Science
Melissa task changed to Development
Melissa task changed to Maintenance
Melissa task changed to Development
edu.bsu.cybersec.core.narrative.InsecurePasswordEvent: Train staff on secure passwords
Expanded: Vani
Expanded: Nancy
Expanded: Melissa
Expanded: Vani
Collapsed: Vani
edu.bsu.cybersec.core.narrative.InputSanitizationEvent: Vani
Final users: 17423; exposure: 0.07671654

6:
CEO is Esteban; workers are Abdullah, Janine, Jerry, 
Expanded: Janine
Collapsed: Janine
Expanded: Jerry
Collapsed: Jerry
Expanded: Abdullah
Collapsed: Abdullah
Jerry task changed to Development
edu.bsu.cybersec.core.narrative.InputSanitizationEvent: Jerry
Abdullah task changed to Maintenance
Janine task changed to Development
edu.bsu.cybersec.core.narrative.DataStolenEvent: Notify our users
Jerry task changed to Development
edu.bsu.cybersec.core.narrative.SecurityConferenceEvent: Jerry
edu.bsu.cybersec.core.narrative.InsecurePasswordEvent: Train staff on secure passwords
edu.bsu.cybersec.core.narrative.ChildAdviceEvent: Computer Science
Jerry task changed to Development
edu.bsu.cybersec.core.narrative.DDOSEvent: Just Wait
Final users: 25900; exposure: 0.096139066

7:
CEO is Nancy; workers are Jerry, Esteban, Janine, 
Expanded: Jerry
Jerry task changed to Development
Collapsed: Jerry
Expanded: Esteban
Collapsed: Esteban
Expanded: Janine
Collapsed: Janine
Expanded: Esteban
Collapsed: Esteban
Expanded: Jerry
Collapsed: Jerry
edu.bsu.cybersec.core.narrative.DDOSEvent: Press Release
Esteban task changed to Development
edu.bsu.cybersec.core.narrative.SecurityConferenceEvent: Jerry
Expanded: Jerry
Collapsed: Jerry
edu.bsu.cybersec.core.narrative.ChildAdviceEvent: Engineering
edu.bsu.cybersec.core.narrative.DataStolenEvent: Notify our users
edu.bsu.cybersec.core.narrative.InsecurePasswordEvent: Train staff on secure passwords
edu.bsu.cybersec.core.narrative.InputSanitizationEvent: Jerry
Final users: 33994; exposure: 0.065729424

8:
CEO is Jerry; workers are Janine, Vani, Esteban, 
Janine task changed to Development
edu.bsu.cybersec.core.narrative.SecurityConferenceEvent: Esteban
edu.bsu.cybersec.core.narrative.DDOSEvent: Just Wait
edu.bsu.cybersec.core.narrative.InsecurePasswordEvent: Train staff on secure passwords
edu.bsu.cybersec.core.narrative.DataStolenEvent: Notify our users
edu.bsu.cybersec.core.narrative.ChildAdviceEvent: Computer Science
edu.bsu.cybersec.core.narrative.InputSanitizationEvent: Esteban
Final users: 38580; exposure: 0.070904195


9-1:
CEO is Jerry; workers are Bruce, Nancy, Janine, 
Expanded: Nancy
Expanded: Bruce
Expanded: Janine
Collapsed: Janine
edu.bsu.cybersec.core.narrative.DataStolenEvent: Notify our users
Expanded: Janine
Collapsed: Janine
Expanded: Nancy
Collapsed: Nancy
Expanded: Bruce
Collapsed: Bruce
edu.bsu.cybersec.core.narrative.DDOSEvent: Just Wait
edu.bsu.cybersec.core.narrative.ChildAdviceEvent: Computer Science
edu.bsu.cybersec.core.narrative.SecurityConferenceEvent: Bruce
edu.bsu.cybersec.core.narrative.InputSanitizationEvent: Nancy
edu.bsu.cybersec.core.narrative.InsecurePasswordEvent: Train staff on secure passwords
Final users: 5493; exposure: 0.012339515

9-2:
CEO is Melissa; workers are Ivar, Esteban, Vani, 
Expanded: Ivar
Ivar task changed to Development
Collapsed: Ivar
Vani task changed to Development
edu.bsu.cybersec.core.narrative.ChildAdviceEvent: Computer Science
edu.bsu.cybersec.core.narrative.SecurityConferenceEvent: Vani
Vani task changed to Development
edu.bsu.cybersec.core.narrative.DDOSEvent: Just Wait
edu.bsu.cybersec.core.narrative.InputSanitizationEvent: Esteban
Vani task changed to Maintenance
edu.bsu.cybersec.core.narrative.InsecurePasswordEvent: Train staff on secure passwords
Final users: 25899; exposure: 0.18964696

10:
CEO is Jerry; workers are Melissa, Bruce, Vani, 
Expanded: Melissa
Expanded: Bruce
Expanded: Vani
Expanded: Bruce
Expanded: Melissa
Expanded: Bruce
Bruce task changed to Development
Expanded: Vani
Vani task changed to Development
Collapsed: Vani
edu.bsu.cybersec.core.narrative.DataStolenEvent: Notify our users
Expanded: Melissa
Collapsed: Melissa
Expanded: Bruce
Collapsed: Bruce
Bruce task changed to Maintenance
edu.bsu.cybersec.core.narrative.ChildAdviceEvent: Computer Science
edu.bsu.cybersec.core.narrative.InputSanitizationEvent: Vani
edu.bsu.cybersec.core.narrative.DDOSEvent: Press Release
edu.bsu.cybersec.core.narrative.SecurityConferenceEvent: Melissa
edu.bsu.cybersec.core.narrative.InsecurePasswordEvent: Train staff on secure passwords
Bruce task changed to Development
Final users: 21280; exposure: 0.094213806

11-1:
CEO is Nancy; workers are Esteban, Jerry, Bruce, 
edu.bsu.cybersec.core.narrative.ChildAdviceEvent: Computer Science
Bruce task changed to Development
edu.bsu.cybersec.core.narrative.DDOSEvent: Press Release
Jerry task changed to Development
edu.bsu.cybersec.core.narrative.DataStolenEvent: Notify our users
edu.bsu.cybersec.core.narrative.InputSanitizationEvent: Esteban
edu.bsu.cybersec.core.narrative.InsecurePasswordEvent: Train staff on secure passwords
edu.bsu.cybersec.core.narrative.SecurityConferenceEvent: Jerry
Esteban task changed to Development
Jerry task changed to Maintenance
Final users: 11928; exposure: 0.15418494

11-2:
CEO is Vani; workers are Janine, Ivar, Abdullah, 
Expanded: Janine
Expanded: Ivar
Expanded: Abdullah
Collapsed: Abdullah
Abdullah task changed to Development
Expanded: Janine
Expanded: Ivar
Expanded: Abdullah
Collapsed: Abdullah
edu.bsu.cybersec.core.narrative.DataStolenEvent: Notify our users
edu.bsu.cybersec.core.narrative.InsecurePasswordEvent: Train staff on secure passwords
Expanded: Ivar
Expanded: Janine
Expanded: Abdullah
Collapsed: Abdullah
edu.bsu.cybersec.core.narrative.InputSanitizationEvent: Janine
edu.bsu.cybersec.core.narrative.ChildAdviceEvent: Computer Science
Abdullah task changed to Development
Abdullah task changed to Maintenance
Janine task changed to Development
edu.bsu.cybersec.core.narrative.DDOSEvent: Press Release
Expanded: Abdullah
Collapsed: Abdullah
edu.bsu.cybersec.core.narrative.SecurityConferenceEvent: Abdullah
Final users: 22768; exposure: 0.096434705

12: 
CEO is Abdullah; workers are Vani, Nancy, Ivar, 
Nancy task changed to Development
edu.bsu.cybersec.core.narrative.DataStolenEvent: Notify our users
Nancy task changed to Maintenance
edu.bsu.cybersec.core.narrative.SecurityConferenceEvent: Vani
Nancy task changed to Development
Ivar task changed to Maintenance
edu.bsu.cybersec.core.narrative.ChildAdviceEvent: Computer Science
Vani task changed to Development
Vani task changed to Maintenance
Ivar task changed to Development
edu.bsu.cybersec.core.narrative.InputSanitizationEvent: Ivar
Ivar task changed to Development
edu.bsu.cybersec.core.narrative.DDOSEvent: Press Release
Set volume to zero
Set volume to zero
edu.bsu.cybersec.core.narrative.InsecurePasswordEvent: Train staff on secure passwords
Final users: 31255; exposure: 0.06462003

13:
CEO is Nancy; workers are Vani, Ivar, Jerry, 
Expanded: Jerry
Jerry task changed to Development
Expanded: Ivar
Ivar task changed to Development
Collapsed: Ivar
edu.bsu.cybersec.core.narrative.ChildAdviceEvent: Computer Science
edu.bsu.cybersec.core.narrative.InputSanitizationEvent: Jerry
Jerry task changed to Development
edu.bsu.cybersec.core.narrative.DDOSEvent: Just Wait
edu.bsu.cybersec.core.narrative.InsecurePasswordEvent: Just change his password
edu.bsu.cybersec.core.narrative.SecurityConferenceEvent: Ivar
edu.bsu.cybersec.core.narrative.InsecurePasswordEvent: Train staff on secure passwords
Vani task changed to Development
Final users: 19108; exposure: 0.15697393
\end{verbatim}


\end{document}
